\documentclass[11pt]{article}
% The shared Artforms notation (a copy of the workspace's notation file).
% Canonical CN-1 printers. Author future notation here; regenerate local copies with migrate.py sync.
% Submission copies are refreshed and verified by migration tooling.
\ifdefined\AFNotationVersion\else
\def\AFNotationVersion{CN-1}
\newcommand{\AFName}[1]{{\mathsf{#1}}}
\newcommand{\AFDepPair}{\Sigma}
\newcommand{\AFDepFun}{\Pi}
\newcommand{\AFFst}{\mathsf{fst}}
\newcommand{\AFSnd}{\mathsf{snd}}
\newcommand{\AFRecognitionProfile}[1][]{{\mathcal P_{\mathrm{rec}\if\relax\detokenize{#1}\relax\else,#1\fi}}}
\newcommand{\AFStructureProfile}{{\mathcal P_{\mathrm{str}}}}
\newcommand{\AFSelectedObservables}{{\mathcal O_{\mathrm{sel}}}}
\newcommand{\AFRecognitionVocabulary}{{\mathcal V_{\mathrm{rec}}}}
\newcommand{\AFFeatureVocabulary}{{\mathcal V_{\mathrm{feat}}}}
\newcommand{\AFDeclSignature}{{\mathfrak S}}
\newcommand{\AFStep}{\mathsf{Step}}
\newcommand{\AFProjector}{\mathsf P}
\newcommand{\AFObservationPartition}{\mathcal B}
\newcommand{\AFCoefficientAutomorphism}{\varphi_R}
\newcommand{\AFCellLoopGroup}{\pi_1^{\mathrm{cell}}}
\newcommand{\AFAddWitness}{\AFName{AddWitness}}
\newcommand{\AFAlignment}{\AFName{Alignment}}
\newcommand{\AFArtformFrame}{\AFName{ArtformFrame}}
\newcommand{\AFAssemble}{\AFName{Assemble}}
\newcommand{\AFAssembleEvents}{\AFName{AssembleEvents}}
\newcommand{\AFAtom}{\AFName{Atom}}
\newcommand{\AFAttachCell}{\AFName{AttachCell}}
\newcommand{\AFBinaryAssemble}{\AFAssemble_2}
\newcommand{\AFCanvas}{\AFName{Canvas}}
\newcommand{\AFCell}{\AFName{Cell}}
\newcommand{\AFCellForm}{\AFName{CellForm}}
% Geometric descriptors and general recognition templates are distinct interfaces.
\newcommand{\AFGeometricForm}{\AFName{GeometricForm}}
\newcommand{\AFRecognitionPattern}{\AFName{RecognitionPattern}}
\newcommand{\AFChromo}{\AFName{Chromo}}
\newcommand{\AFChromoForm}{\AFName{ChromoForm}}
\newcommand{\AFChromoState}{\AFName{ChromoState}}
\newcommand{\AFClock}{\AFName{Clock}}
\newcommand{\AFClockFiltration}{\AFName{ClockFiltration}}
\newcommand{\AFCoeffIntWit}{\AFName{CoeffIntWit}}
\newcommand{\AFCoefficientProfile}{\AFName{CoefficientProfile}}
\newcommand{\AFCoherentRealizer}{\AFName{CoherentRealizer}}
\newcommand{\AFCompatibleBatch}{\AFName{CompatibleBatch}}
\newcommand{\AFConcurrencyWitness}{\AFName{ConcurrencyWitness}}
\newcommand{\AFConnection}{\AFName{Connection}}
\newcommand{\AFContent}{\AFName{Content}}
\newcommand{\AFCoordIntWit}{\AFName{CoordIntWit}}
\newcommand{\AFCountWitness}{\AFName{CountWitness}}
\newcommand{\AFCoxeterBall}{\AFName{CoxeterBall}}
\newcommand{\AFCubicalGrid}{\AFName{CubicalGrid}}
\newcommand{\AFCuboidView}{\AFName{CuboidView}}
\newcommand{\AFCuboidViewL}{\AFCuboidView_L}
\newcommand{\AFCuboidViewR}{\AFCuboidView_R}
\newcommand{\AFCyclarchy}{\AFName{Cyclarchy}}
\newcommand{\AFCycleTime}{\AFName{CycleTime}}
\newcommand{\AFCylinder}{\AFName{Cylinder}}
\newcommand{\AFDefectFusionWitness}{\AFName{DefectFusionWitness}}
\newcommand{\AFEmpty}{\AFName{Empty}}
\newcommand{\AFEventTapestry}{\AFName{EventTapestry}}
\newcommand{\AFFeatureState}{\AFName{FeatureState}}
\newcommand{\AFField}{\AFName{Field}}
\newcommand{\AFFinCode}{\AFName{FinCode}}
\newcommand{\AFFold}{\AFName{Fold}}
\newcommand{\AFForm}{\AFName{Form}}
\newcommand{\AFGaugeAlign}{\AFName{GaugeAlign}}
\newcommand{\AFGaugeFrame}{\AFName{GaugeFrame}}
\newcommand{\AFGaugeObsIntWit}{\AFName{GaugeObsIntWit}}
\newcommand{\AFGaugeSafe}{\AFName{GaugeSafe}}
\newcommand{\AFGaugeSafeRecIntWit}{\AFName{GaugeSafeRecIntWit}}
\newcommand{\AFGlue}{\AFName{Glue}}
\newcommand{\AFGlueBoundary}{\AFName{GlueBoundary}}
\newcommand{\AFGridView}{\AFName{GridView}}
\newcommand{\AFGroupSource}{\AFName{GroupSource}}
\newcommand{\AFHistory}{\AFName{History}}
\newcommand{\AFIncidenceCanvas}{\AFName{IncidenceCanvas}}
\newcommand{\AFIndex}{\AFName{Index}}
\newcommand{\AFInstantiation}{\AFName{Instantiation}}
\newcommand{\AFInstrument}{\AFName{Instrument}}
\newcommand{\AFIntWit}{\AFName{IntWit}}
\newcommand{\AFInterpret}{\AFName{Interpret}}
\newcommand{\AFLawSpec}{\AFName{LawSpec}}
\newcommand{\AFLinearTime}{\AFName{LinearTime}}
\newcommand{\AFLinkTable}{\AFName{LinkTable}}
\newcommand{\AFLocalization}{\AFName{Localization}}
\newcommand{\AFLocalize}{\AFName{Localize}}
\newcommand{\AFMatch}{\AFName{Match}}
\newcommand{\AFMatchIntWit}{\AFName{MatchIntWit}}
\newcommand{\AFMaxwellState}{\AFName{MaxwellState}}
\newcommand{\AFMulWitness}{\AFName{MulWitness}}
\newcommand{\AFNatForm}{\AFName{NatForm}}
\newcommand{\AFNoClaim}{\AFName{NoClaim}}
\newcommand{\AFNormalize}{\AFName{Normalize}}
\newcommand{\AFObsIntWit}{\AFName{ObsIntWit}}
\newcommand{\AFObservable}{\AFName{Observable}}
\newcommand{\AFObserve}{\AFName{Observe}}
\newcommand{\AFObserverProfile}{\AFName{ObserverProfile}}
\newcommand{\AFOccurrence}{\AFName{Occurrence}}
\newcommand{\AFOrb}{\AFName{Orb}}
\newcommand{\AFOrderedAssemble}{\AFAssemble^{\mathrm{ord}}}
\newcommand{\AFOrderedForm}{\AFName{OrderedForm}}
\newcommand{\AFOrientCanvas}{\AFName{OrientCanvas}}
\newcommand{\AFOrientedCellForm}{\AFName{OrientedCellForm}}
\newcommand{\AFPatchWitness}{\AFName{PatchWitness}}
\newcommand{\AFPathComplex}{\AFName{PathComplex}}
\newcommand{\AFPerfectAlign}{\AFName{PerfectAlign}}
\newcommand{\AFPhaseGridPair}{\AFName{PhaseGridPair}}
\newcommand{\AFPlacement}{\AFName{Placement}}
\newcommand{\AFPointTime}{\AFName{PointTime}}
\newcommand{\AFPointedInt}{\AFName{PointedInt}}
\newcommand{\AFPointedTimeForm}{\AFName{PointedTimeForm}}
\newcommand{\AFPolygonComplex}{\AFName{PolygonComplex}}
\newcommand{\AFPos}{\AFName{Pos}}
\newcommand{\AFPoseOrbit}{\AFName{PoseOrbit}}
\newcommand{\AFPresentation}{\AFName{Presentation}}
\newcommand{\AFProcessProfile}{\AFName{ProcessProfile}}
\newcommand{\AFProductRod}{\AFName{ProductRod}}
\newcommand{\AFProtocol}{\AFName{Protocol}}
\newcommand{\AFProtoform}{\AFName{Protoform}}
% Legacy printer retained for manuscripts not yet migrated; never aliases a shape type.
\newcommand{\AFQuotient}{\AFName{Quotient}}
\newcommand{\AFRatioForm}{\AFName{RatioForm}}
\newcommand{\AFRealization}{\AFName{Realization}}
\newcommand{\AFRealize}{\AFName{Realize}}
\newcommand{\AFRealizer}{\AFName{Realizer}}
\newcommand{\AFRectTapestry}{\AFName{RectTapestry}}
\newcommand{\AFRefine}{\AFName{Refine}}
\newcommand{\AFRefineTime}{\AFName{RefineTime}}
\newcommand{\AFRendering}{\AFName{Rendering}}
\newcommand{\AFReplicate}{\AFName{Replicate}}
\newcommand{\AFRewriteApplication}{\AFName{RewriteApplication}}
\newcommand{\AFRod}{\AFName{Rod}}
\newcommand{\AFRootPatch}{\AFName{RootPatch}}
\newcommand{\AFRun}{\AFName{Run}}
\newcommand{\AFRunContract}{\AFName{RunContract}}
\newcommand{\AFSampleClock}{\AFName{SampleClock}}
\newcommand{\AFSampledPicture}{\AFName{SampledPicture}}
\newcommand{\AFScheduler}{\AFName{Scheduler}}
\newcommand{\AFSchedulerPolicy}{\AFName{SchedulerPolicy}}
\newcommand{\AFScript}{\AFName{Script}}
\newcommand{\AFSegmentedRod}{\AFName{SegmentedRod}}
\newcommand{\AFSegmentedRodL}{\AFSegmentedRod_L}
\newcommand{\AFSegmentedRodR}{\AFSegmentedRod_R}
\newcommand{\AFSem}{\AFName{Sem}}
\newcommand{\AFSplit}{\AFName{Split}}
\newcommand{\AFSquareGrid}{\AFName{SquareGrid}}
\newcommand{\AFState}{\AFName{State}}
\newcommand{\AFSubstitute}{\AFName{Substitute}}
\newcommand{\AFSubstitutionStage}{\AFName{SubstitutionStage}}
\newcommand{\AFSweep}{\AFName{Sweep}}
\newcommand{\AFTally}{\AFName{Tally}}
\newcommand{\AFTapestryForm}{\AFName{TapestryForm}}
\newcommand{\AFTemporalCode}{\AFName{TemporalCode}}
\newcommand{\AFTimeForm}{\AFName{TimeForm}}
\newcommand{\AFToken}{\AFName{Token}}
\newcommand{\AFTorsor}{\AFName{Torsor}}
\newcommand{\AFTorus}{\AFName{Torus}}
\newcommand{\AFTraceTapestry}{\AFName{TraceTapestry}}
\newcommand{\AFTrajectory}{\AFName{Trajectory}}
\newcommand{\AFTriangularGrid}{\AFName{TriangularGrid}}
\newcommand{\AFTypedAtom}{\AFName{TypedAtom}}
\newcommand{\AFValuedAtom}{\AFName{ValuedAtom}}
\newcommand{\AFWheel}{\AFName{Wheel}}
\newcommand{\AFWitness}{\AFName{Witness}}

% Temporal companion records: derived interfaces, not primitive carriers.
\newcommand{\AFOccurrenceHistory}{\AFName{OccurrenceHistory}}
\newcommand{\AFTickRealization}{\AFName{TickRealization}}
\newcommand{\AFDurationRecord}{\AFName{DurationRecord}}
\newcommand{\AFEventAssociation}{\AFName{EventAssociation}}
\newcommand{\AFCalibrationRecord}{\AFName{CalibrationRecord}}
\newcommand{\AFElapsed}{\AFName{Elapsed}}

% Spatial measurement: derived records and evaluators.
\newcommand{\AFMeasurandSpec}{\AFName{MeasurandSpec}}
\newcommand{\AFQuantityValue}{\AFName{QuantityValue}}
\newcommand{\AFMetricContentRecord}{\AFName{MetricContentRecord}}
\newcommand{\AFSpatialAssociation}{\AFName{SpatialAssociation}}
\newcommand{\AFSpatialCalibrationRecord}{\AFName{SpatialCalibrationRecord}}
\newcommand{\AFSpatialMeasurementWitness}{\AFName{SpatialMeasurementWitness}}
\newcommand{\AFSpatialRead}{\AFName{SpatialRead}}
\fi



\usepackage[a4paper,margin=1in]{geometry}
\usepackage[T1]{fontenc}
\usepackage[utf8]{inputenc}
\usepackage{lmodern}
\usepackage{microtype}
\usepackage{amsmath,amssymb,amsthm,mathtools}
\usepackage{booktabs,tabularx,array,longtable}
\usepackage{enumitem}
\usepackage[dvipsnames]{xcolor}
\usepackage{tikz}
\usetikzlibrary{arrows.meta,positioning,calc}
\usepackage[font=small,labelfont=bf]{caption}
\usepackage[colorlinks=true,linkcolor=MidnightBlue,citecolor=MidnightBlue,urlcolor=MidnightBlue]{hyperref}
\usepackage[nameinlink,noabbrev]{cleveref}

\numberwithin{equation}{section}
\setcounter{tocdepth}{1}

\theoremstyle{definition}
\newtheorem{definition}{Definition}[section]
\newtheorem{example}[definition]{Example}
\newtheorem{principle}[definition]{Principle}
\newtheorem{observation}[definition]{Observation}
\theoremstyle{plain}
\newtheorem{theorem}[definition]{Theorem}
\newtheorem{proposition}[definition]{Proposition}
\newtheorem{lemma}[definition]{Lemma}
\newtheorem{corollary}[definition]{Corollary}
\theoremstyle{remark}
\newtheorem{remark}[definition]{Remark}

\crefname{definition}{Definition}{Definitions}
\crefname{example}{Example}{Examples}
\crefname{principle}{Principle}{Principles}
\crefname{observation}{Observation}{Observations}
\crefname{theorem}{Theorem}{Theorems}
\crefname{proposition}{Proposition}{Propositions}
\crefname{lemma}{Lemma}{Lemmas}
\crefname{corollary}{Corollary}{Corollaries}

\newcommand{\RATT}{\mathsf{RATT}_{\mathrm c}}
\newcommand{\Type}{\mathsf{type}}
\renewcommand{\Form}{\AFForm}
\newcommand{\Hist}{\AFHistory}
\newcommand{\Pres}{\AFPresentation}
\newcommand{\Script}{\AFScript}
\newcommand{\State}{\AFState}
\newcommand{\Graphic}{\mathsf{Graphic}}
\newcommand{\Event}{\mathsf{Event}}
\newcommand{\Origin}{\mathsf{Origin}}
\newcommand{\Stamped}{\mathsf{Stamped}}
\newcommand{\Role}{\mathsf{Role}}
\newcommand{\Frm}{\mathsf F}
\newcommand{\Prs}{\mathsf P}
\newcommand{\Wit}{\mathsf W}
\newcommand{\Scr}{\mathsf S}
\newcommand{\Ren}{\mathsf R}
\newcommand{\FORM}{\textsf{FORM}}
\newcommand{\PRES}{\textsf{PRESENTATION}}
\newcommand{\WIT}{\textsf{WITNESS}}
\newcommand{\SCR}{\textsf{SCRIPT}}
\newcommand{\REN}{\textsf{RENDERING}}
\newcommand{\Gd}{\mathbb G}
\newcommand{\eps}{\varepsilon}
\newcommand{\eff}{\mathsf{eff}}
\newcommand{\valid}{\mathsf{valid}}
\newcommand{\atoms}{\mathsf{At}}
\newcommand{\role}{\mathsf{role}}
\newcommand{\cost}{\mathsf{cost}}
\newcommand{\erase}{\mathsf{erase}}
\newcommand{\events}{\mathsf{events}}
\newcommand{\nf}{\mathsf{nf}}
\newcommand{\supp}{\mathsf{supp}}
\newcommand{\unit}{\mathbf e}
\newcommand{\absent}{\boldsymbol\infty}
\newcommand{\jtype}{\mathrel{\mathsf{type}}}
\newcommand{\jdeq}{\equiv}
\newcommand{\ctx}{\mathsf{ctx}}
\newcommand{\wf}{\mathsf{wf}}
\newcommand{\old}{\mathsf{old}}
\newcommand{\new}{\mathsf{new}}
\newcommand{\cross}{\mathsf{Cross}}
\newcommand{\sem}[1]{\left\lbrack\!\left\lbrack #1\right\rbrack\!\right\rbrack}
\newcommand{\fst}{\AFFst}
\newcommand{\snd}{\AFSnd}
\newcommand{\id}{\mathsf{id}}
\newcommand{\fin}{\mathsf{Fin}}

\tikzset{
  role/.style={draw=MidnightBlue!65,rounded corners=3pt,fill=MidnightBlue!5,
    minimum width=20mm,minimum height=9mm,align=center},
  arr/.style={-{Latex[length=2mm]},thick,MidnightBlue!75},
  gatearr/.style={-{Latex[length=2mm]},thick,dashed,BrickRed!75},
  forbidarr/.style={-{Latex[length=2mm]},thick,densely dotted,Gray!70}
}

\hypersetup{
  pdftitle={Five Roles for Artforms},
  pdfauthor={Sergey Krivov},
  pdfsubject={A role-stratified type theory with traceable gates},
  pdfkeywords={artforms, Artform Calculus, dependent type theory, graded effects, noninterference, provenance}
}

\title{\textbf{Five Roles for Artforms}\\[4pt]
\large A Role-Stratified Type Theory with Traceable Gates}
\author{Sergey Krivov}
\date{October 2026}

\begin{document}
\maketitle

\begin{abstract}
The textbook \emph{Artform Calculus} sorts the data of every construction
into five roles, \FORM, \PRES, \WIT, \SCR{} and \REN, and states a policy for
them: what has been built, presented, certified or done may be used
downstream; a witness or a script reaches the frame only through a gate that
records what was assumed; and nothing may depend on a reading.  This paper
makes that discipline precise and tests the claim that five is the right
number.

The policy is generated by four levels of a construction: the carrier, its
carried structure, the instance, and the display.  It is the only policy that
meets four requirements the book states, and it is path closed, which is
exactly the condition under which checking each direct dependency suffices.
For constructions whose dependencies are graded \emph{free}, \emph{gated} or
\emph{absent} we prove the firewall, mediation by recorded gates, frame
stability and presentation independence, for chains of dependencies and for
computed values.  The policy tells every two roles apart except \WIT{} and
\SCR, which share the instance level and must: a certificate may mention a
move, and a guarded move takes its certificate.  Their rules tell them apart.
A history is a list of applications of scripts, and a certificate only admits
an application; counting certificates as values, or applications by their
endpoints, changes the counts the book makes.  A state is a legal choice of
values, which is witness data, so the five roles classify every datum.  These
results and the book's examples are checked in Lean~4 with Mathlib.

The second part gives \(\RATT\), a universe-free dependent calculus whose
judgements carry the role.  Every rule computes an effect in the min--max
algebra of the three grades, and a registered gate consumes a stamped
credential and contracts to a private node that records its event.  We give
paper proofs of principal effects, substitution, subject reduction, strong
normalization, confluence, canonical forms and decidable conversion; of
static--dynamic gate adequacy; and of module theorems: rendering-only and
fresh-witness-family extensions are conservative at the frame's grade-free
types, and new evidence may enable a construction but never silently.  The
book's running examples are written in the calculus.
\end{abstract}

\tableofcontents

\section{Introduction}\label{sec:intro}

The textbook \emph{Artform Calculus} \cite{Book} teaches one discipline from
its first chapter: every datum of a construction has a role, and the role
decides what may depend on it.  There are five roles.  A \FORM{} retains a
carrier and its construction history.  A \PRES{} retains a selected order,
seams, coordinates, cuts or sectors.  A \WIT{} retains evidence that a
comparison or a move is allowed.  A \SCR{} is a total action on the legal
states of a presentation.  A \REN{} is one chosen visible reading, and no
rule consumes it \cite[\S1.9]{Book}.  The book's chess example puts all five
on one board: the squares with their adjacency are form, the addressing
\(a1\) to \(h8\) is presentation, the certificate that a move is legal is a
witness, the move is a script, and the printed diagram and an engine's
evaluation are renderings.

The book claims more than a vocabulary.  It says that the list is closed,
that a sixth role would make a different calculus, and that one rule carries
most of the weight: nothing may depend on a reading.  This paper asks what it
takes for those claims to be solid, and proves what it can.  Four questions
organize the answer.
\begin{enumerate}[label=(\arabic*),leftmargin=2.2em]
\item \emph{Is the policy forced?}  The book states the policy in prose.
Four requirements that the book states determine it uniquely
(\cref{thm:forced}).
\item \emph{Does checking each dependency suffice?}  A policy is useful only
if a construction can be checked one dependency at a time.  That is so
exactly for path-closed policies, and the book's policy is path closed
(\cref{thm:path-closed,thm:local-global}).
\item \emph{Do the guarantees hold?}  Nothing outside the rendering depends
on a reading; every dependency the policy gates crosses a gate whose record
reaches the result; the frame changes only on the record; and a form never
depends on its addressing except on the record.  Each holds for chains of
dependencies and for computed values (\cref{sec:constructions}).
\item \emph{Are there exactly five roles?}  The policy tells nine of the ten
pairs of roles apart.  It cannot tell \WIT{} from \SCR, and should not; their
rules tell them apart, and the book's counts of histories depend on the
difference (\cref{sec:histories}).  A state, which the book declines to call
a role, is witness data, so every datum has one of the five roles
(\cref{sec:states}).
\end{enumerate}

\Cref{sec:roles,sec:constructions,sec:histories,sec:examples} answer these
questions for the role discipline alone.  Every theorem there is checked in
Lean~4 with Mathlib \cite{Lean4,Mathlib}, in the library \texttt{ArtformRoles}
of the repository that accompanies the book \cite{Repo}; \cref{tab:lean} lists
the Lean names.  The proofs use only the axioms \texttt{propext},
\texttt{Classical.choice} and \texttt{Quot.sound}.

The book keeps the rules of inference, which are Martin-L\"of's
\cite{MartinLof}, and the roles as two disciplines: the rules check that each
value has its type, and only the roles say what a value may be used for
\cite[\S1.11]{Book}.  \Cref{sec:core,sec:reduction,sec:metatheory,sec:adequacy,sec:modules,sec:erasure}
put the two into one judgement.  The calculus \(\RATT\), a role-stratified
artform type theory, is a universe-free dependent type theory whose
judgement
\(\AFDeclSignature;\Gamma\vdash_{\rho}t:A\triangleright\eps\) carries the role
\(\rho\) of its result and the effect \(\eps\) of every declaration it uses.
An effect is a vector of grades, and its paths are exactly the dependency
chains of \cref{sec:constructions} (\cref{lem:path-characterization}).  A
registered gate is the only rule that may introduce a gated dependency, and
its contraction leaves a node that records the event which authorized it.
The metatheory there is proved on paper; its mechanization is the open task
of \cref{sec:status}.  \Cref{sec:artforms} writes the book's examples in the
calculus.

\subsection{Contributions}

The paper contributes:
\begin{enumerate}[label=(\roman*),leftmargin=2.5em]
\item the book's five roles with the book's readings, and a decision that
makes them classify every datum: a state is witness data
(\cref{sec:states});
\item the generation of the policy by four levels, its uniqueness under the
book's four requirements, and its path closure
(\cref{thm:levels,thm:forced,thm:path-closed});
\item constructions with graded dependencies, on which checking each direct
dependency suffices exactly for path-closed policies, and the firewall,
recorded mediation, frame stability and presentation independence for
chains and for values (\cref{sec:constructions});
\item the separation of all ten pairs of roles: nine by the policy, and
\WIT{} from \SCR{} by the count of histories (\cref{thm:separation});
\item Lean proofs of all of the above, and of the book's chess, realizer and
walker examples (\cref{sec:examples,tab:lean});
\item the calculus \(\RATT\) with principal min--max effects, graded dependent
functions, stamped credentials, registered gates and private record nodes,
and paper proofs of substitution, preservation, strong normalization,
confluence, canonical forms, decidable conversion, path mediation and
static--dynamic gate adequacy;
\item module theorems that distinguish rendering vocabulary, fresh witness
families and new inhabitants of old evidence families, at an explicit
grade-free boundary; and
\item the book's frames, states, scripts, interpretation witnesses and
renderings written in \(\RATT\).
\end{enumerate}

\subsection{Relation to an earlier manuscript}

This paper revises the manuscript \emph{Traceable Gates for Artforms}
(September 2026), whose calculus is the second part here.  There the roles
were written \(\mathsf O,\mathsf S,\mathsf W,\mathsf T,\mathsf R\) and read as
constructed objects, installed structure, evidence, transformations and
renderings.  Here they carry the book's names and readings, and the policy is
the same matrix.  Three things are new: the first part, which justifies the
five roles and is checked in Lean; the decision that a state is witness data;
and the word \emph{record} for what a gate contraction leaves, since the book
uses \emph{receipt} for a witness of a measured feature
\cite[Prop.~3.43]{Book}.

\section{The five roles}\label{sec:roles}

\subsection{Roles}\label{sec:role-readings}

\Cref{tab:roles} gives the five roles with the book's readings
\cite[\S1.9]{Book}.  A role is a job that a datum does at a stage of a
construction, not a kind of substance: the base cut of a wheel is a witness
while it is offered to the constructor that points the wheel, and
presentation once the pointed wheel carries it.  The role of a datum is the
role of the type at which it is judged at that stage, so the question a role
answers is never what a thing is, but what it is doing now.

\begin{table}[ht]
\centering
\small
\begin{tabularx}{\textwidth}{@{}lc>{\raggedright\arraybackslash}X>{\raggedright\arraybackslash}X@{}}
\toprule
Role & & What it retains & Typical examples\\
\midrule
\FORM & \(\Frm\) & carrier and construction history & cells, tags, order,
incidence, orientation\\
\PRES & \(\Prs\) & selected order, seams, coordinates, cuts, sectors & rod,
grid view, wheel, coordinate system\\
\WIT & \(\Wit\) & evidence that a comparison or move is allowed, and the
choices it certifies & alignment, contact, a measurement certificate, a
supplied policy, a state\\
\SCR & \(\Scr\) & total action on the legal states of a presentation &
flatten, reflection, update, forgetful passage, readout\\
\REN & \(\Ren\) & one chosen visible reading, consumed by no rule & drawing,
table, sound, a value read off\\
\bottomrule
\end{tabularx}
\caption{The five roles.  The readings are the book's; the witness row adds
the choices a witness certifies, among them a state (\cref{sec:states}).}
\label{tab:roles}
\end{table}

\begin{definition}[Roles]\label{def:roles}
The roles are
\(
  \Role=\{\Frm,\Prs,\Wit,\Scr,\Ren\}
\),
read as in \cref{tab:roles}.  Throughout, every role metavariable
\(\sigma,\tau,\kappa,\rho\), with or without an index, ranges over \(\Role\).
\end{definition}

A \emph{frame} is a typed specification of a carrier, its carried structure,
the space of legal states, the allowed observations and the allowed moves; an
\emph{instance} is one legal state inhabiting a frame \cite[Def.~1.13]{Book}.
The carrier and its construction are form, the carried structure is
presentation, the moves are scripts, and an observation is a script whose
values are renderings.  The frame does not change while a run proceeds; the
instance does.

\subsection{States}\label{sec:states}

The book calls a state one side of the frame/instance split rather than a
job, and declines to make it a sixth role.  It also says which role a state's
data have: when a role must be named for a state, it is a witness
\cite[\S1.9]{Book}.  Earlier editions said instead to name the presentation
that the state inhabits; the reasons for the present answer follow.

A state is a choice of values at the positions of a frame, made legal by
evidence.  In the book's formation rule the state space of a presentation
\(A\) is the tagged choice of a table with the evidence that it is legal,
\(\AFDepPair_{s}\mathsf{Legal}_A(s)\), and the legality is counted as a
property: different proofs of the same legality claim do not make different
states \cite[\S1.9]{Book}.  Choices and evidence are the witness role's data.
The book counts a supplied policy and a chosen batch as witnesses, and says
of its process laboratory that evidence and choice are one role.  So a state
is witness data.  It is not carried structure, since a law changes it while
the frame stays fixed; and it is not a script, since it does not act.  The
book's connection, a witness family that became a state, keeps its role when
it is gathered into one state \cite[\S5.9]{Book}, and a derived state such as
the squares one side bears on is witness data that a rule may consult
(\cref{sec:chess}).  The rule concerns the states of a frame, whose positions
are fixed.  A process that grows its carrier, as the book's growth does,
has forms for its stages, and each step of growth is a gate into a form.

\begin{principle}[Every datum has one of the five roles]\label{pr:closed}
In a construction every datum is form, presentation, witness, script or
rendering.  The book's four useful words that are not roles have these
homes:
\begin{description}[leftmargin=2.2cm,style=sameline]
\item[state] witness data: a legal choice of values;
\item[policy] witness data while it is supplied, presentation once a gate
installs it (\cref{sec:realizer});
\item[reduct] the result of a script that forgets part of a vocabulary;
\item[observable] a script whose values are renderings.
\end{description}
\end{principle}

\subsection{Grades and the policy}\label{sec:policy}

A dependency of one datum on another is used in one of three ways: in the
ordinary way, through a gate that records what was assumed, or not at all.
Let
\[
  \Gd=\{0<1<\absent\},
  \qquad
  p\oplus q=\min(p,q),
  \qquad
  p\otimes q=\max(p,q).
\]
Grade \(0\) is ordinary use, grade \(1\) is use through a gate, and
\(\absent\) is no use.  Two dependencies in sequence compose with \(\otimes\);
of two alternative routes, the better counts, by \(\oplus\).  Then
\((\Gd,\oplus,\absent,\otimes,0)\) is an idempotent commutative semiring.

The book states its policy in prose \cite[\S1.9]{Book}.  Whatever has been
built, presented, certified or done may be used freely downstream; a witness
or a script may feed back into a presentation or a form only through a gate
that records what was assumed; and nothing may depend on a reading, by any
route through a third role.  \Cref{tab:policy} is the policy as a matrix
\(D:\Role\times\Role\to\Gd\).  Its entry \(D(\sigma,\rho)\) is the least
grade at which a datum of role \(\rho\) may depend on a datum of role
\(\sigma\).  It is the matrix of the earlier manuscripts, under the book's
names.  The book's prose names every gated entry, including the gate from
presentation to form, which its editions before October 2026 left implicit
(\cref{sec:levels}).

\begin{table}[ht]
\centering
\small
\begin{tabular}{@{}lccccc@{}}
\toprule
source / result & \(\Frm\) & \(\Prs\) & \(\Wit\) & \(\Scr\) & \(\Ren\)\\
\midrule
\(\Frm\) & 0 & 0 & 0 & 0 & 0\\
\(\Prs\) & 1 & 0 & 0 & 0 & 0\\
\(\Wit\) & 1 & 1 & 0 & 0 & 0\\
\(\Scr\) & 1 & 1 & 0 & 0 & 0\\
\(\Ren\) & \(\absent\) & \(\absent\) & \(\absent\) & \(\absent\) & 0\\
\bottomrule
\end{tabular}
\caption{The dependency policy \(D\).  Rows are the roles depended on,
columns the roles that depend.  A demand of \(\absent\) means that only an
absent dependency is accepted.}
\label{tab:policy}
\end{table}

\subsection{Four levels}\label{sec:levels}

\begin{definition}[Levels]\label{def:levels}
The \emph{level} of a role is
\[
 \ell(\Frm)=0,\qquad
 \ell(\Prs)=1,\qquad
 \ell(\Wit)=\ell(\Scr)=2,\qquad
 \ell(\Ren)=3,
\]
read as the frame's carrier, the frame's carried structure, the instance,
and the display.  For any map \(\lambda:\Role\to\mathbb N\) let
\[
 D_\lambda(\sigma,\rho)=
 \begin{cases}
 \absent&\text{if }\sigma=\Ren\ne\rho,\\
 0&\text{otherwise, if }\lambda(\sigma)\le\lambda(\rho),\\
 1&\text{otherwise.}
 \end{cases}
\]
\end{definition}

\begin{theorem}[The policy has four levels]\label{thm:levels}
\(D=D_\ell\).  A dependency is free from a lower or equal level, gated from a
higher one, and absent out of a reading.
\end{theorem}

The book gives four reasons for the policy, and each is a requirement on a
level map \(\lambda\).
\begin{enumerate}[label=(R\arabic*),leftmargin=2.8em]
\item \emph{Nothing may depend on a reading, and a reading may be taken of
anything}: \(\lambda(\sigma)<\lambda(\Ren)\) for every \(\sigma\ne\Ren\).
\item \emph{The frame does not change while a run proceeds}; a witness or a
script reaches a form or a presentation only through a gate:
\(\lambda(\Frm),\lambda(\Prs)<\lambda(\Wit),\lambda(\Scr)\).
\item \emph{No addressing is part of the board}: a form does not depend on
how it is addressed, so \(\lambda(\Frm)<\lambda(\Prs)\).
\item \emph{A certificate may mention a move, and a guarded move takes its
certificate}: the certificate that a capture is legal names the capture, and
a guarded move is total on its admitted input, including the needed evidence
\cite[\S1.9]{Book}.  Neither direction is feedback, so
\(\lambda(\Wit)=\lambda(\Scr)\).
\end{enumerate}

\begin{theorem}[The policy is forced]\label{thm:forced}
The levels \(\ell\) satisfy \textup{(R1)--(R4)}, and every map
\(\lambda:\Role\to\mathbb N\) satisfying \textup{(R1)--(R4)} has
\(D_\lambda=D\).
\end{theorem}

\begin{proof}
The requirements fix every comparison between the values of \(\lambda\) on
two roles, and \(D_\lambda\) depends only on those comparisons and on the
firewall clause.
\end{proof}

\subsection{Path closure}

A policy is \emph{path closed} when
\begin{equation}\label{eq:path-closure}
  D(\sigma,\rho)\le D(\sigma,\tau)\otimes D(\tau,\rho)
  \qquad\text{for all }\sigma,\tau,\rho.
\end{equation}
The book's key idea says that no route through a third role recovers a
forbidden dependency; path closure is that sentence for every pair of steps.

\begin{theorem}[Path closure]\label{thm:path-closed}
The policy \(D\) is path closed, and \(D(\sigma,\sigma)=0\) for every role.
\end{theorem}

For example, a script may freely produce evidence, but a dependency of a
presentation on that evidence still has demand \(1\); and a witness built
from a reading would carry the reading's demand \(\absent\) into everything
built from the witness.  \Cref{thm:local-global} shows that path closure is
exactly what makes one-step checking sound.

\subsection{What the policy tells apart}\label{sec:policy-separation}

Say that the policy \emph{separates} two roles \(\sigma\) and \(\tau\) when
some role \(\kappa\) has \(D(\sigma,\kappa)\ne D(\tau,\kappa)\) or
\(D(\kappa,\sigma)\ne D(\kappa,\tau)\): some datum may depend on one of them
in a way it may not depend on the other, or one of them may depend on
something in a way the other may not.

\begin{theorem}[Policy separation]\label{thm:policy-separation}
The policy separates \(\sigma\) and \(\tau\) if and only if
\(\sigma\ne\tau\) and \(\{\sigma,\tau\}\ne\{\Wit,\Scr\}\).  The rows of
\(\Wit\) and \(\Scr\) are equal, and so are their columns.
\end{theorem}

\Cref{tab:separation} names, for each separated pair, a role on which the two
differ and the guarantee of \cref{sec:constructions} that the difference
carries.  Merging the two roles of a pair would either forbid a dependency
the book makes or permit one the guarantee excludes.

\begin{table}[ht]
\centering
\small
\begin{tabularx}{\textwidth}{@{}l>{\raggedright\arraybackslash}p{.36\textwidth}X@{}}
\toprule
Pair & Where the demands differ & Guarantee\\
\midrule
\(\Frm,\Prs\) & \(D(\Prs,\Frm)=1\), \(D(\Frm,\Frm)=0\) & presentation
independence (\cref{thm:form-stable})\\
\(\Frm,\Wit\);\ \(\Frm,\Scr\) & \(D(\Wit,\Frm)=D(\Scr,\Frm)=1\),
\(D(\Frm,\Frm)=0\) & frame stability for forms (\cref{thm:frame-stable})\\
\(\Prs,\Wit\);\ \(\Prs,\Scr\) & \(D(\Wit,\Prs)=D(\Scr,\Prs)=1\),
\(D(\Prs,\Prs)=0\) & frame stability for presentations
(\cref{thm:frame-stable})\\
\(\sigma,\Ren\) for \(\sigma\ne\Ren\) & \(D(\Ren,\sigma)=\absent\),
\(D(\sigma,\sigma)=0\) & the firewall (\cref{cor:firewall,thm:firewall-values})\\
\(\Wit,\Scr\) & nowhere & the count of histories (\cref{thm:separation})\\
\bottomrule
\end{tabularx}
\caption{How each pair of roles is told apart.}
\label{tab:separation}
\end{table}

\section{Constructions}\label{sec:constructions}

\subsection{Graded dependencies}

\begin{definition}[Construction]\label{def:construction}
A \emph{construction} on a set \(V\) of items gives each item \(v\) a role
\(\role(v)\) and each pair of items \(u,v\) a grade \(c(u,v)\in\Gd\): the
grade at which \(v\) depends directly on \(u\), with \(\absent\) for no
dependency.  It is \emph{valid} for a policy \(D\) when
\[
  D(\role(u),\role(v))\le c(u,v)\qquad\text{for all }u,v\in V.
\]
For a grade \(g\), a \emph{\(g\)-chain} from \(u\) to \(v\) is a sequence
\(u=v_0,v_1,\ldots,v_k=v\) with \(c(v_{i-1},v_i)\le g\) for every \(i\); the
empty chain has \(k=0\).  A \(0\)-chain uses ordinary dependencies only; a
\(1\)-chain is a \emph{live} chain.
\end{definition}

The derivations of \cref{sec:core} give constructions: the items are the
declarations and variables a derivation uses and its result, and the grades
are the coefficients of its effect equations.  \Cref{lem:path-characterization}
shows that the effect of a derivation is the best grade over its chains.

\begin{theorem}[Local to global]\label{thm:local-global}
Let \(D\) be a policy with \(D(\sigma,\sigma)=0\) for every role.
\begin{enumerate}[label=(\roman*),leftmargin=2.5em]
\item If \(D\) is path closed, then in every construction valid for \(D\),
a \(g\)-chain from \(u\) to \(v\) exists only if \(D(\role(u),\role(v))\le g\).
\item Conversely, if \textup{(i)} holds for every construction on three items,
then \(D\) is path closed.
\end{enumerate}
\end{theorem}

\begin{proof}
(i) By induction on the chain: the empty chain needs
\(D(\sigma,\sigma)=0\), and a chain extended by one dependency needs one
instance of \eqref{eq:path-closure}.  (ii) For roles \(\sigma,\tau,\rho\),
give three items these roles, let the first two dependencies have grades
\(D(\sigma,\tau)\) and \(D(\tau,\rho)\), and let every other grade be
\(\absent\).  The construction is valid and has a
\(D(\sigma,\tau)\otimes D(\tau,\rho)\)-chain from the first item to the third.
\end{proof}

From now on a construction is valid for the book's policy \(D\).

\begin{corollary}[Firewall]\label{cor:firewall}
No live chain leads from a rendering to a datum of another role.
\end{corollary}

\begin{corollary}[Mediation]\label{cor:mediation}
If \(D(\role(u),\role(v))=1\), no chain of ordinary dependencies leads from
\(u\) to \(v\).
\end{corollary}

\begin{theorem}[Records]\label{thm:records}
If \(D(\role(u),\role(v))=1\) and a live chain leads from \(u\) to \(v\),
then the chain crosses a gate: there are items \(w,w'\) with
\(c(w,w')=1\), a live chain from \(u\) to \(w\), and a live chain from
\(w'\) to \(v\).  The record the gate leaves at \(w'\) reaches \(v\).
\end{theorem}

\begin{proof}
A live chain either has only ordinary dependencies, which
\cref{cor:mediation} excludes, or contains a dependency of grade \(1\);
split it there.
\end{proof}

\subsection{Values}\label{sec:values}

The statements so far are about chains.  The guarantees the book wants are
about values: what a reading changes, and what a run may change.

\begin{definition}[Evaluation]\label{def:evaluation}
An \emph{evaluation} of a construction with values in a set \(X\) gives each
item \(v\) a map \(f_v:X^V\to X\) that depends only on the values of the
items \(u\) with \(c(u,v)\le1\).  A valuation \(x\in X^V\) is
\emph{consistent} with it when \(x_v=f_v(x)\) for every \(v\).  A
construction is \emph{acyclic} when some rank \(r:V\to\mathbb N\) has
\(r(u)<r(v)\) whenever \(c(u,v)\le1\).
\end{definition}

\begin{lemma}[Determination]\label{lem:determination}
Let a construction be acyclic, let \(U\subseteq V\), and let two evaluations
agree on the items of \(U\).  If \(x\) and \(y\) are consistent with the two
evaluations and \(x_u=y_u\) whenever \(c(u,v)\le1\) for some \(v\in U\) and
\(u\notin U\), then \(x\) and \(y\) agree on \(U\).
\end{lemma}

\begin{proof}
By induction on the rank of the item.
\end{proof}

\begin{theorem}[Firewall for values]\label{thm:firewall-values}
Replacing every renderer of an acyclic construction, and nothing else,
changes no value outside the rendering: if two evaluations agree on every
item that is not a rendering, consistent valuations agree there.
\end{theorem}

\begin{theorem}[Frame stability]\label{thm:frame-stable}
Let two evaluations of an acyclic construction agree on its forms and
presentations, and let consistent valuations \(x\) and \(y\) agree on the
source of every gate into a form or a presentation.  Then \(x\) and \(y\)
agree on every form and presentation.  The frame changes only through
recorded gates.
\end{theorem}

\begin{theorem}[Presentation independence]\label{thm:form-stable}
Let two evaluations of an acyclic construction agree on its forms, and let
consistent valuations agree on the source of every gate into a form.  Then
they agree on every form.  In particular a form depends on how it is
addressed only through recorded gates.
\end{theorem}

Each of the three is \cref{lem:determination} for one set \(U\): the items
that are not renderings, the frame's items, and the forms.  Validity says
that every live dependency entering \(U\) from outside it is a gate, or, for
the firewall, that there is none.

\subsection{Derived is not displayed}

Two values may be computed from one state by explicit computation and still
have different roles.  The book's example is a chess position: the squares
one side bears on form a table with a value on every square, and a rule may
consult it; an engine's evaluation is one number, and no rule of chess
consults it \cite[\S1.9]{Book}.  In a construction the difference is the role
of the item, which is the role of the type at which it is judged: the codomain
decides, not the amount of computation.  The first table is witness data, a
derived state; the evaluation is a rendering, and \cref{cor:firewall} forbids
every live chain from it to a certificate (\cref{sec:chess}).

The same test settles the book's other pairs: curvature, a value on every
face, is a derived state, while a holonomy read off one loop is a reading; a
conserved charge is reported, while the divergence clause on every site is a
constraint between states.

\section{Witnesses and scripts}\label{sec:histories}

\subsection{What the policy cannot do}

\Cref{thm:policy-separation} leaves one pair: witnesses and scripts have the
same row and the same column of the policy.  Requirement (R4) explains why,
and the book needs both directions.  Witnesses depend on scripts: the
certificate of a capture names the capture, the reversibility of a shear is a
property of the shear, and the evidence of a realizer says that every batch
it selects keeps the invariants.  Scripts depend on witnesses: a guarded move
takes its certificate, and the book's spatial reading of a measurement
witness is a script built from that witness.  Making either direction a gate
would demand a record where the book has none.

Their rules tell them apart.  A script acts: it has a source and a target,
composes at the presentation in between, and has an action on legal states.
A witness admits or chooses: it is consumed by a guard or a gate, and when it
admits an application, which certificate it is does not matter.  The book
makes the difference count \cite[\S7.4]{Book}.

\subsection{Histories}

A law on a finite set \(X\) of labels gives, for labels \(y\) and \(z\), the
finite set \(\mathsf{App}(y,z)\) of its applications from \(y\) to \(z\),
which are script data, and for each application \(a\) the finite set
\(\mathsf{Ev}(a)\) of its certificates, which are witness data.  An
application is \emph{admitted} when it has a certificate.  A history of
\(n\) steps is a list of \(n\) admitted applications, each starting where the
previous one ended.

The book weights histories to obtain its tallies.  With weight \(1\) a
history counts once; with nonnegative weights it gives the probabilities of a
stochastic realizer; with the coefficients of a coherent law it gives the
quantum tallies, and interference is cancellation between two histories with
one endpoint.  In each case the book counts applications and not proofs: if
each proof of an application's legality counted as a value, one move would be
two histories and its coefficient would be counted twice
\cite[\S7.4]{Book}.

\begin{definition}[Three counts]\label{def:counts}
Fix a start \(x_0\in X\).  For \(n\ge0\) and \(z\in X\) let
\(h_0(z)=p_0(z)=t_0(z)=[z=x_0]\) and
\[
\begin{aligned}
 h_{n+1}(z)&=\sum_{y\in X}\ \sum_{a\in\mathsf{App}(y,z)}
   [\mathsf{Ev}(a)\ne\varnothing]\;h_n(y),\\
 p_{n+1}(z)&=\sum_{y\in X}\ \sum_{a\in\mathsf{App}(y,z)}
   |\mathsf{Ev}(a)|\;p_n(y),\\
 t_{n+1}(z)&=\sum_{y\in X}
   \bigl[\exists a\in\mathsf{App}(y,z).\ \mathsf{Ev}(a)\ne\varnothing\bigr]\;t_n(y).
\end{aligned}
\]
The \emph{histories} \(h_n(z)\) count lists of admitted applications, the
book's count.  The \emph{proof lists} \(p_n(z)\) count every certificate as a
value.  The \emph{trajectories} \(t_n(z)\) count sequences of labels joined by
admitted applications, so that applications count only by their endpoints.
\end{definition}

\begin{proposition}\label{prop:counts-agree}
If every application has at most one certificate, then \(p_n=h_n\).  If
between any two labels at most one application is admitted, then
\(h_n=t_n\).
\end{proposition}

\begin{example}[Two corners]\label{ex:two-corners}
Let \(X=\{0,1\}\).  From either corner the token may step ahead or behind,
and both steps reach the other corner.  Stepping ahead has two certificates
and stepping behind has one.  After one step from \(0\),
\[
 p_1(1)=3,\qquad h_1(1)=2,\qquad t_1(1)=1.
\]
\end{example}

\subsection{Five roles}

\begin{theorem}[Separation]\label{thm:separation}
Every two distinct roles are told apart.
\begin{enumerate}[label=(\roman*),leftmargin=2.5em]
\item If \(\{\sigma,\tau\}\ne\{\Wit,\Scr\}\), the policy separates \(\sigma\)
and \(\tau\), and \cref{tab:separation} names the guarantee the difference
carries.
\item Witnesses and scripts are separated by the count of histories: on the
law of \cref{ex:two-corners}, counting applications and certificates alike as
values gives \(3\), counting both alike as properties gives \(1\), and the
book's count, which counts applications as values and certificates as
properties, gives \(2\).
\end{enumerate}
\end{theorem}

\begin{proof}
(i) is \cref{thm:policy-separation}.  (ii) is \cref{ex:two-corners}: a single
role for applications and certificates would have to count both as \(p\)
counts them, or both as \(t\) counts them.
\end{proof}

With \cref{pr:closed} this answers the fourth question.  The policy has four
levels, and (R1)--(R4) force them.  Each level holds one role, except the
instance level, which holds the two roles that act and that admit or choose.
A sixth role would need either a fifth level, and so a fifth requirement the
book does not state, or a third job at the instance level that is neither
acting nor admitting nor choosing.

\section{The book's examples}\label{sec:examples}

\subsection{One capture}\label{sec:chess}

The book's \S1.9 follows one capture, \(N{\times}e5\), with the five roles.
\Cref{tab:chess} lists its items as a construction.  The board's squares with
their adjacency and the step shape of each piece are form; the addressing is
presentation; the positions before and after the capture, the squares White
bears on, and the certificate that the capture is legal are witness data; the
capture is a script; the game score is a form that keeps the moves; the
engine's evaluation and the printed diagram are renderings.  The score keeps
the capture through a gate, and every other dependency is ordinary.

\begin{table}[ht]
\centering
\small
\begin{tabularx}{\textwidth}{@{}llX@{}}
\toprule
Item & Role & Depends on\\
\midrule
board & \(\Frm\) & ---\\
addressing & \(\Prs\) & board\\
position & \(\Wit\) & board, addressing\\
squares White bears on & \(\Wit\) & board, position\\
certificate of \(N{\times}e5\) & \(\Wit\) & position, squares White bears
on, the capture\\
capture \(N{\times}e5\) & \(\Scr\) & ---\\
position after & \(\Wit\) & position, capture, certificate\\
game score & \(\Frm\) & capture, through a gate\\
evaluation & \(\Ren\) & position\\
diagram & \(\Ren\) & position, addressing\\
\bottomrule
\end{tabularx}
\caption{One capture as a construction.}
\label{tab:chess}
\end{table}

\begin{proposition}\label{prop:chess}
The capture of \cref{tab:chess} is valid.  It becomes invalid if the
certificate depends on the evaluation, or if the score depends on the capture
without a gate.  No live chain leads from the evaluation to the certificate
or to the position after the capture.
\end{proposition}

\subsection{A realizer}\label{sec:realizer}

A law says which moves are allowed and chooses none; a realizer is a policy
together with evidence that every batch it selects keeps the law's invariants
\cite[Def.~5.10]{Book}.  The policy is one datum at two stages.  While it is
supplied it is witness data; once a realizer carries it, it is presentation.
As a construction: the law is a form, the supplied policy a witness that
depends on the law, and the realizer a presentation that depends on the law
and, through a gate, on the supplied policy.  The construction is valid, no
chain of ordinary dependencies leads from the supplied policy to the
realizer, and a live chain does, across the installation gate
(\cref{cor:mediation,thm:records}).

\subsection{The walker}\label{sec:walker}

The book's walker moves a token round the four corners of a square, one step
ahead or behind at every tick, and each step has one certificate: that the
token is at the corner it leaves \cite[Def.~5.9]{Book}.  From \(v_0\) there are
four histories of two ticks; two end at \(v_0\) and two at \(v_2\).  Each
step has one certificate and leads to a corner of its own, so by
\cref{prop:counts-agree} the three counts of \cref{def:counts} agree on the
square for every number of ticks.  They come apart only when a step has two
certificates, or two steps share an endpoint, as on the two corners of
\cref{ex:two-corners}.

\begin{table}[ht]
\centering
\small
\begin{tabularx}{\textwidth}{@{}>{\raggedright\arraybackslash}p{.43\textwidth}>{\raggedright\arraybackslash}X@{}}
\toprule
Result & Lean name\\
\midrule
semiring of grades & \texttt{Grade.semiring\_laws}\\
\cref{thm:levels} & \texttt{demand\_eq\_demandOf\_level}\\
\cref{thm:forced} & \texttt{requirements\_level}, \texttt{demandOf\_eq\_demand}\\
\cref{thm:path-closed} & \texttt{demand\_pathClosed}, \texttt{demand\_self}\\
\cref{thm:policy-separation} & \texttt{policySeparates\_iff},
\texttt{witness\_script\_same\_policy}\\
\cref{thm:local-global} & \texttt{Construction.Valid.reach},
\texttt{Construction.pathClosed\_of\_reach}\\
\cref{cor:firewall,cor:mediation} & \texttt{Construction.no\_reach\_of\_rendering},
\texttt{Construction.no\_free\_reach}\\
\cref{thm:records} & \texttt{Construction.crosses\_gate}\\
\cref{lem:determination} & \texttt{Construction.Evaluation.eq\_of\_agree}\\
\cref{thm:firewall-values,thm:frame-stable,thm:form-stable} &
\texttt{Construction.firewall}, \texttt{Construction.frame\_stable},
\texttt{Construction.form\_stable}\\
\cref{prop:counts-agree} & \texttt{Law.proofs\_eq\_histories},
\texttt{Law.histories\_eq\_trajectories}\\
\cref{ex:two-corners}, \cref{thm:separation}(ii) & \texttt{twoCorners\_counts}\\
\cref{prop:chess} & \texttt{chess\_valid}, \texttt{chess\_readback\_invalid},
\texttt{chess\_silent\_score\_invalid}, \texttt{chess\_evaluation\_terminal}\\
\Cref{sec:realizer} & \texttt{realizer\_valid},
\texttt{realizer\_installation\_recorded}\\
\Cref{sec:walker} & \texttt{walker\_two\_ticks}, \texttt{walker\_counts\_agree}\\
\bottomrule
\end{tabularx}
\caption{The results of \cref{sec:roles,sec:constructions,sec:histories,sec:examples}
in the Lean library \texttt{ArtformRoles}, namespace
\texttt{Artform.Roles}.}
\label{tab:lean}
\end{table}

\section{The core calculus}\label{sec:core}

This section and the five after it give \(\RATT\), the dependent calculus
whose judgements carry the roles.  Its policy is the matrix \(D\) of
\cref{tab:policy}, and its grades are those of \cref{sec:policy}.  The roles
\(\Frm\) and \(\Prs\), whose data make up the frame, are called
\emph{protected}: a gated dependency can only end in them.

\subsection{Effects}\label{sec:effects}

For a finite set of atoms \(Z\), an \emph{effect} is a vector
\(
  \eps:Z\to\Gd
\)
with default value \(\absent\).  Define
\[
\begin{aligned}
 (\eps\oplus\eta)(z)&=\eps(z)\oplus\eta(z),\\
 (q\otimes\eps)(z)&=q\otimes\eps(z),\\
 \unit_z(z)&=0,
 &\unit_z(y)&=\absent\quad(y\ne z).
\end{aligned}
\]
Restriction deleting \(x\) is written \(\eps\setminus x\).  Closed literals and
meta-level finite indices have no declaration atom and therefore introduce no
dependency.  Let \(\atoms(\AFDeclSignature,\Gamma)\) contain the public declarations
of \(\AFDeclSignature\) and the variables of \(\Gamma\), each with its declared role.

\begin{definition}[Policy-valid effect]\label{def:valid-effect}
An effect \(\eps\) at result role \(\rho\) is policy valid, written
\(\valid_\rho(\eps)\), when
\[
  D(\role(z),\rho)\le\eps(z)
  \qquad(z\in\atoms(\AFDeclSignature,\Gamma)).
\]
Thus demand \(1\) requires every live path from \(z\) to contain a gate, and
demand \(\absent\) requires that no live path remain.
\end{definition}

A value computed from an artform has the role of the type at which it is
judged: the codomain decides, and there is no role-free escape hatch for
derived data.

\subsection{The universe-free boundary}\label{sec:boundary}

The common Artforms notation uses \(\AFDeclSignature\) for a declaration
signature and \(\AFDepPair\), \(\AFDepFun\) for dependent type formers.  The
role and grade annotations in this paper are part of their checked
interfaces: pairs remain homogeneous in role, and \(\Type\) remains a
judgmental sort.  Upright sans-serif names share the corpus spelling; no
ambient datum thereby becomes an internal declaration.

\(\RATT\) has dependent functions and pairs but no internal universe.  Types
are classified by a meta-level judgment rather than packaged as terms.  This
choice removes the standard counterexample
\(
  \AFDepPair_{X:\mathsf{Type}}X
\), in which a fresh type and one constructor form a new inhabitant of an old
type.  Role-indexed universes are an important extension, but including them
before their elimination rules are proved noninterfering would weaken the main
theorem.

The core also excludes arbitrary user rewrite rules, quotient computation,
and extensions of script equality.  Strictly positive inductive families and
certified recursion may be added through a host-kernel interface, but are not
needed to state or prove the gate metatheory.  Forced algebraic-value
extension from the Pattern Geometry manuscript remains a separate staged
problem \cite{PatternGeometry}.

\subsection{Syntactic categories}

The calculus is an LF-style dependent type theory: \(\Type\) is a judgmental
sort, not an internal term.  A role is an index of a declaration or judgment,
not part of the type itself.  Here and in every rule below,
\(\sigma,\rho\in\Role\).  We therefore write \(x:_{\sigma}A\) for a
declaration at source role \(\sigma\), and \(\vdash_{\rho}\) for a formation
or typing judgment at result role \(\rho\).  The subscripted colon is also
used in binders and telescopes; \(\Rightarrow_{\rho}\) marks the result role
of a family or gate declaration.  The unindexed colon in
\(\vdash_{\rho}t:A\) retains its ordinary meaning: \(t\) has type \(A\).
Grades decorate dependent function spaces.

The five-element set \(\Role\) is the global range from which role indices
are drawn, but the role attached to a particular well-formed type is not an
arbitrary choice from that set.  We call \(\rho\) an \emph{admissible outer
role} of \(A\) when a formation judgment
\(\AFDeclSignature;\Gamma\vdash_{\rho}A\;\jtype\) is derivable.  The formation rules
below make this outer role unique.  Parameters and proper subexpressions of
\(A\) may nevertheless carry different roles.

\begin{align*}
A,B ::={}&
  a(\vec t)
  \mid \Event
  \mid \Origin_A(e,t)
  \mid \Stamped_\sigma(A)
  \mid \AFDepFun^q(x:_{\sigma}A).B
  \mid \AFDepPair(x:_{\rho}A).B,
\\
t,u ::={}&
  x\mid c\mid
  \lambda^q(x:_{\sigma}A).t\mid t\,u\mid
  (t,u)\mid\fst(t)\mid\snd(t)\\
 &\mid \mathsf{stamp}(t,e,o)
  \mid \mathsf{gate}_g(\vec d;s)
  \mid \mathsf{node}_g(\vec d;t,e).
\end{align*}
Here \(q\in\Gd\), \(a\) is a declared type-family symbol, and \(c\) is a
declared public constant.  The last form is internal syntax: users cannot
write \(\mathsf{node}_g\), but reduction can produce it.  Dependent sums are
homogeneous in role.  Cross-role packaging is provided only by the sealed
\(\Stamped\) former and consumed only by a gate.

An expression is \emph{public} when it contains no private
\(\mathsf{node}_g\) form.  Write \(\mathsf{gradefree}(A)\) when every
\(\AFDepFun\)-annotation occurring anywhere in \(A\) is \(0\).  This is a
syntactic observation boundary, not a claim that grade-zero functions are
first order: it excludes only types that can retain a forbidden dependency
behind a positive or absent latent function edge.

For exact attestation, let \(\mathsf{wdata}(A)\) be generated by
\[
\frac{
 a:(\Delta)\Rightarrow_{\Wit}\Type
}{
 \mathsf{wdata}(a(\vec t))
}
\qquad
\frac{
 \mathsf{wdata}(A)\quad\mathsf{wdata}(B)
}{
 \mathsf{wdata}\bigl(\AFDepPair(x:_{\Wit}A).B\bigr)
}.
\tag{WitnessData}
\]
Thus witness data consists of witness-family applications and finite
homogeneous dependent sums of them.  Event, origin, stamp, function,
rendering, and non-witness heads are excluded.

A context is
\[
  \Gamma ::= \cdot\mid\Gamma,x:_{\rho}A.
\]
The principal judgments are
\[
\begin{array}{rcl}
\AFDeclSignature\;\wf,&&
\AFDeclSignature\vdash\Gamma\;\ctx,\\[2pt]
\AFDeclSignature;\Gamma\vdash_{\rho} A\;\jtype\triangleright\eps,&&
\AFDeclSignature;\Gamma\vdash_{\rho} t:A\triangleright\eps.
\end{array}
\]
Effects record free variables and public declaration atoms.  Every displayed
formation or typing rule has the side condition that its conclusion effect is
policy valid.  We write a superscript \(\checkmark\) on a premise when that
side condition has already been discharged.

\subsection{Signatures}

A signature is a finite ordered list of the following entries.

\begin{description}[leftmargin=3.2cm,style=nextline]
\item[Family]
\(
  a:(\Delta)\Rightarrow_{\rho}\Type
\), where \(\Delta\) is an already well-formed telescope.
\item[Primitive]
\(
  c:_{\rho}A
\), where \(A\) is closed and policy valid.  Every \(\AFDepFun\)-grade in the
type of a primitive is \(0\), and its ultimate result head is a declared
family symbol.  Thus a primitive cannot merely assert that it mediates an
argument, introduce an opaque pair, or manufacture an unowned credential.
\item[Certificate bundle]
For closed \(\vdash_{\sigma}A\;\jtype\) with \(\sigma\ne\Ren\), no nonzero
\(\AFDepFun\)-grade, and ordinary result head, a bundle
\[
 d=(a_d,e_d,o_d):A
\]
introduces fresh public constants
\[
 a_d:_{\sigma}A,
 \qquad e_d:_{\Wit}\Event,
 \qquad o_d:_{\Wit}\Origin_A(e_d,a_d).
\]
The three names have the same module owner.
\item[Exact attestation]
For a closed public term
\[
 \vdash_{\Wit}t:A,
 \qquad
 \mathsf{gradefree}(A),
 \qquad
 \mathsf{wdata}(A),
\]
an entry \(q@t=(e_q,o_q):A\) introduces only
\[
 e_q:_{\Wit}\Event,
 \qquad
 o_q:_{\Wit}\Origin_A(e_q,t).
\]
The two fresh names have the same module owner.  The payload is the exact
already constructed term \(t\), not a new constant.
\item[Gate]
A gate entry
\[
 g:\Delta;\Stamped_\sigma(A)\Longrightarrow_{\rho} B
 \quad[\mathsf{node}_g]
\]
requires \(\rho\in\{\Frm,\Prs\}\), \(D(\sigma,\rho)\le1\), and:
\begin{enumerate}[label=(G\arabic*),leftmargin=3em]
\item \(\Delta\) is an ordinary parameter telescope whose declaration roles
have demand \(0\) into \(\rho\);
\item \(\Delta\vdash_{\sigma}A\;\jtype\) is derivable;
\item \(\Delta\vdash_{\rho}B\;\jtype\) is derivable, and \(B\) does not depend on
the stamped payload, and contains no nonzero \(\AFDepFun\)-grade; and
\item \(\mathsf{node}_g\) is a fresh private constructor with ordinary
parameter edges and grade-\(1\) payload and event edges.
\end{enumerate}
The public gate has the graded type
\[
  \AFDepFun^0(\vec x:\Delta).\,
  \AFDepFun^1(s:_{\Wit}\Stamped_\sigma(A)).\,B.
\]
\end{description}

All names are globally qualified by their module.  There is no declaration
form for a custom reduction rule, no public declaration of a positive
\(\AFDepFun\)-grade, no positive grade in a gate's result type, no attestation of
an open or non-witness-data term, and no eliminator from \(\Stamped\).  These
negative clauses are part of signature
well-formedness, not optional module advice.

\subsection{Type formation and effects}\label{sec:type-rules}

The family rule uses ordinary edges:
\[
\frac{
 a:(x_1:_{\rho_1}A_1,\ldots,x_n:_{\rho_n}A_n)\Rightarrow_{\rho}\Type\in\AFDeclSignature
 \quad
 \AFDeclSignature;\Gamma\vdash_{\rho_i} t_i:A_i[\vec t_{<i}/\vec x_{<i}]
 \triangleright\eps_i
}{
 \AFDeclSignature;\Gamma\vdash_{\rho} a(\vec t)\;\jtype
 \triangleright
 \unit_a\oplus\bigoplus_i\eps_i
}.
\tag{Ty-Fam}
\]
The atom \(a\) makes a newly declared family visible even when it has no
arguments.  The built-in event type has no public declaration atom:
\[
\frac{}{\AFDeclSignature;\Gamma\vdash_{\Wit}\Event\;\jtype
\triangleright\absent}.
\tag{Ty-Event}
\]
Here and below \(\absent\) also denotes the everywhere-\(\absent\) vector.

Origin and stamping are type formers at the witness role:
\[
\frac{
 \AFDeclSignature;\Gamma\vdash_{\sigma} A\;\jtype\triangleright\alpha
 \quad
 \AFDeclSignature;\Gamma\vdash_{\Wit} e:\Event\triangleright\eta
 \quad
 \AFDeclSignature;\Gamma\vdash_{\sigma} t:A\triangleright\theta
}{
 \AFDeclSignature;\Gamma\vdash_{\Wit}\Origin_A(e,t)\;\jtype
 \triangleright\alpha\oplus\eta\oplus\theta
}
\tag{Ty-Origin}
\]
and
\[
\frac{
 \AFDeclSignature;\Gamma\vdash_{\sigma} A\;\jtype\triangleright\alpha
 \quad \sigma\ne\Ren
}{
 \AFDeclSignature;\Gamma\vdash_{\Wit}\Stamped_\sigma(A)\;\jtype
 \triangleright\alpha
}.
\tag{Ty-Stamp}
\]

The grade on a dependent function is both a reduction-stable lower bound on
its operational argument use and the edge by which its domain type enters the
function type.  The explicit premise \(D(\sigma,\rho)\le q\) checks the role
of the binder itself, even when its domain has no public family atom (as for
\(\Event\)):
\[
\frac{
 \AFDeclSignature;\Gamma\vdash_{\sigma} A\;\jtype\triangleright\alpha
 \quad
 \AFDeclSignature;\Gamma,x:_{\sigma}A\vdash_{\rho} B\;\jtype\triangleright\beta
 \quad q\in\Gd
 \quad D(\sigma,\rho)\le q
}{
 \AFDeclSignature;\Gamma\vdash_{\rho}\AFDepFun^q(x:_{\sigma}A).B\;\jtype
 \triangleright(q\otimes\alpha)\oplus(\beta\setminus x)
}.
\tag{Ty-Pi}
\]
The premise for \(B\) separately checks any occurrence of \(x\) in the
codomain.  In particular, a gate result type cannot hide a witness-dependent
index; gate results depend only on their ordinary telescope.

Dependent sums are homogeneous and have ordinary component edges:
\[
\frac{
 \AFDeclSignature;\Gamma\vdash_{\rho} A\;\jtype\triangleright\alpha
 \quad
 \AFDeclSignature;\Gamma,x:_{\rho}A\vdash_{\rho} B\;\jtype\triangleright\beta
}{
 \AFDeclSignature;\Gamma\vdash_{\rho}\AFDepPair(x:_{\rho}A).B\;\jtype
 \triangleright\alpha\oplus(\beta\setminus x)
}.
\tag{Ty-Sigma}
\]

\begin{lemma}[Role functionality]\label{lem:role-functionality}
For a fixed well-formed signature \(\AFDeclSignature\) and context \(\Gamma\), if
\[
 \AFDeclSignature;\Gamma\vdash_{\rho}A\;\jtype\triangleright\alpha
 \qquad\text{and}\qquad
 \AFDeclSignature;\Gamma\vdash_{\rho'}A\;\jtype\triangleright\alpha',
\]
then \(\rho=\rho'\).  Hence every well-formed type has exactly one
admissible outer role.  In a declaration \(x:_{\sigma}A\) or a term judgment
\(\vdash_{\rho}t:A\), the displayed role must be this unique outer role of
\(A\); it is not independently chosen from \(\Role\).
\end{lemma}

\begin{proof}
By inversion on the two formation derivations.  For a family application,
the unique declaration of its head symbol fixes the result role.  The
\(\Event\), \(\Origin\), and \(\Stamped\) formers have the fixed result role
\(\Wit\).  A dependent function has the outer role of its codomain, so the
claim follows inductively from the two derivations of that codomain.  A
dependent sum is homogeneous and its displayed component role is also its
outer role.  These are all type formers.  Policy acceptance can reject a raw
formation derivation but cannot change its role.
\end{proof}

\subsection{Terms and principal effects}\label{sec:term-rules}

Variables and public constants introduce unit effects:
\[
\frac{x:_{\rho}A\in\Gamma}
{\AFDeclSignature;\Gamma\vdash_{\rho} x:A\triangleright\unit_x}
\tag{Var}
\]
\[
\frac{c:_{\rho}A\in\AFDeclSignature}
{\AFDeclSignature;\Gamma\vdash_{\rho} c:A\triangleright\unit_c}.
\tag{Const}
\]

For abstraction, the declared grade is a lower bound on the principal path
effect of the bound parameter in the body.
\[
\frac{
 \AFDeclSignature;\Gamma\vdash_{\sigma} A\;\jtype
 \quad
 \AFDeclSignature;\Gamma,x:_{\sigma}A\vdash_{\rho} t:B\triangleright\theta
 \quad q\le\theta(x)
}{
 \AFDeclSignature;\Gamma\vdash_{\rho}\lambda^q(x:_{\sigma}A).t:
 \AFDepFun^q(x:_{\sigma}A).B
 \triangleright\theta\setminus x
}.
\tag{Lam}
\]
Application composes the argument paths with the latent grade:
\[
\frac{
 \AFDeclSignature;\Gamma\vdash_{\rho} f:\AFDepFun^q(x:_{\sigma}A).B\triangleright\phi
 \quad
 \AFDeclSignature;\Gamma\vdash_{\sigma} u:A\triangleright\eta
}{
 \AFDeclSignature;\Gamma\vdash_{\rho} f\,u:B[u/x]
 \triangleright\phi\oplus(q\otimes\eta)
}.
\tag{App}
\]
For source elaboration, an omitted grade is chosen maximally and exactly as
\(q=\theta(x)\).  An explicit smaller grade is a conservative promise: it may
report more dependence than the current body exhibits, but never less.  This
slack is required by subject reduction.  For example, reducing
\(\lambda^0x.\fst(c,x)\) to \(\lambda^0x.c\) must not change its function
type.  The binder-policy premise in \textsc{Ty-Pi} prevents such slack from
turning a demand-\(1\) or demand-\(\absent\) role crossing into a grade-zero
interface.

Thus a maximally elaborated unused parameter receives grade \(\absent\), an
ordinary use receives grade \(0\), and a parameter used only beneath gates
receives grade \(1\).
Every term judgment also requires a prior accepted formation derivation for
its displayed type; this convention is made explicit in \cref{app:rules}.

Pairing and projections are ordinary:
\[
\frac{
 \AFDeclSignature;\Gamma\vdash_{\rho} t:A\triangleright\theta
 \quad
 \AFDeclSignature;\Gamma\vdash_{\rho} u:B[t/x]\triangleright\eta
}{
 \AFDeclSignature;\Gamma\vdash_{\rho}(t,u):\AFDepPair(x:_{\rho}A).B
 \triangleright\theta\oplus\eta
},
\tag{Pair}
\]
\[
\frac{\AFDeclSignature;\Gamma\vdash_{\rho} p:\AFDepPair(x:_{\rho}A).B\triangleright\pi}
{\AFDeclSignature;\Gamma\vdash_{\rho}\fst(p):A\triangleright\pi}
\tag{Fst}
\]
\[
\frac{\AFDeclSignature;\Gamma\vdash_{\rho} p:\AFDepPair(x:_{\rho}A).B\triangleright\pi}
{\AFDeclSignature;\Gamma\vdash_{\rho}\snd(p):B[\fst(p)/x]\triangleright\pi}.
\tag{Snd}
\]

A stamp is a sealed triple.  It has introduction but no user-accessible
projection or recursor:
\[
\frac{
 \AFDeclSignature;\Gamma\vdash_{\sigma} t:A\triangleright\theta
 \quad
 \AFDeclSignature;\Gamma\vdash_{\Wit} e:\Event\triangleright\eta
 \quad
 \AFDeclSignature;\Gamma\vdash_{\Wit} o:\Origin_A(e,t)\triangleright\omega
}{
 \AFDeclSignature;\Gamma\vdash_{\Wit}\mathsf{stamp}(t,e,o):
 \Stamped_\sigma(A)
 \triangleright\theta\oplus\eta\oplus\omega
}.
\tag{Stamp}
\]

For a registered gate \(g:\Delta;\Stamped_\sigma(A)\Rightarrow_{\rho} B\),
write \(\delta\) for the ordinary parameter effect and \(\eta\) for the
stamp effect.  Its public and private rules are
\[
\frac{
 \AFDeclSignature;\Gamma\vdash\vec d:\Delta\triangleright\delta
 \quad
 \AFDeclSignature;\Gamma\vdash_{\Wit} s:\Stamped_\sigma(A[\vec d])
 \triangleright\eta
}{
 \AFDeclSignature;\Gamma\vdash_{\rho}\mathsf{gate}_g(\vec d;s):B[\vec d]
 \triangleright\unit_g\oplus\delta\oplus(1\otimes\eta)
}
\tag{Gate}
\]
and
\[
\frac{
 \AFDeclSignature;\Gamma\vdash\vec d:\Delta\triangleright\delta
 \quad
 \AFDeclSignature;\Gamma\vdash_{\sigma} t:A[\vec d]\triangleright\theta
 \quad
 \AFDeclSignature;\Gamma\vdash_{\Wit} e:\Event\triangleright\eta
}{
 \AFDeclSignature;\Gamma\vdash_{\rho}\mathsf{node}_g(\vec d;t,e):B[\vec d]
 \triangleright\unit_g\oplus\delta\oplus
  (1\otimes(\theta\oplus\eta))
}.
\tag{Node}
\]
Only the reduction relation may invoke \textsc{Node}.  The owner atom remains
\(g\); the private constructor does not create a second public capability.

Finally, conversion changes only the displayed type:
\[
\frac{
 \AFDeclSignature;\Gamma\vdash_{\rho} t:A\triangleright\theta
 \quad A\jdeq_{\rho} B
}{
 \AFDeclSignature;\Gamma\vdash_{\rho} t:B\triangleright\theta}.
\tag{Conv}
\]

\begin{example}[A higher-order gate wrapper]\label{ex:wrapper}
Fix a registered gate
\(
  g:\Delta;\Stamped_\Wit(A)\Rightarrow_{\Frm} B
\)
and closed ordinary parameters \(\vec d:\Delta\).  In the body
\[
  \mathsf{gate}_g(\vec d;c)
\]
the bound credential \(c\) has effect component \(1\).  Consequently
\[
  G\;=\;\lambda^1(c:_{\Wit}\Stamped_\Wit(A)).
       \mathsf{gate}_g(\vec d;c)
\]
has type
\(
  \AFDepFun^1(c:_{\Wit}\Stamped_\Wit(A)).B
\).
For a certificate bundle \((a_d,e_d,o_d)\), application reduces as
\[
 G\,\mathsf{stamp}(a_d,e_d,o_d)
 \longrightarrow
 \mathsf{gate}_g(\vec d;\mathsf{stamp}(a_d,e_d,o_d))
 \longrightarrow
 \mathsf{node}_g(\vec d;a_d,e_d).
\]
Thus the latent grade permits higher-order abstraction without replacing the
record-producing gate that justifies it.
\end{example}

All side conditions, simultaneous telescope rules, and context/signature
formation rules are collected without abbreviation in
\cref{app:rules}.

\section{Reduction and definitional equality}\label{sec:reduction}

\subsection{One-step reduction}

Capture-avoiding substitution is written \(t[u/x]\).  The primitive redexes
are
\begin{align}
 (\lambda^q(x:_{\sigma}A).t)\,u
   &\longrightarrow t[u/x],
   \tag{$\beta_q$}\label{red:beta}\\
 \fst(t,u)&\longrightarrow t,
   \tag{$\pi_1$}\label{red:fst}\\
 \snd(t,u)&\longrightarrow u,
   \tag{$\pi_2$}\label{red:snd}\\
 \mathsf{gate}_g(\vec d;\mathsf{stamp}(t,e,o))
   &\longrightarrow \mathsf{node}_g(\vec d;t,e).
   \tag{record}\label{red:gate}
\end{align}
The reduction relation is the compatible closure of these axioms under every
type and term former.  In particular, it reduces under binders, in dependent
indices, in each stamp field, in every gate parameter, and in every retained
node field.  The full one-hole-context grammar is in
\cref{app:reduction}.

The record rule deliberately drops the origin proof but retains both the
payload and event.  A private node has no head reduction.  Hence a consumed
stamp cannot later erase its event field while the corresponding node
survives.  The retained event may itself reduce or be transformed by a later
substitution.

Definitional equality \(t\jdeq_{\rho} u:A\) and type equality
\(A\jdeq_{\rho} B\) are the least typed, reflexive, symmetric, transitive, and
congruent relations containing \(\longrightarrow\).  There are no
eta-rules and no user equations in the core.

\subsection{Normal forms and event support}

A term is normal when it contains no instance of
\cref{red:beta,red:fst,red:snd,red:gate}.  For a normal term \(n\), define its
event support structurally by
\[
\begin{aligned}
 \events(\mathsf{node}_g(\vec d;t,e))
   &=\{e\}\cup\events(\vec d)\cup\events(t),\\
 \events(k(\vec t))&=\bigcup_i\events(t_i)
 \quad\text{for every other constructor }k.
\end{aligned}
\]
Events inside an unused argument erased by beta reduction do not occur in the
normal support.  Events inserted by a contracted gate do occur, because a node
has no event-erasing rule.

\begin{definition}[Residual record trace]\label{def:cross}
Instrument each instance of \cref{red:gate} with a fresh metatheoretic label
\(\ell\) on the node it creates.  Labels are ignored by typing and reduction,
are copied when substitution duplicates a node, and disappear only when the
whole labelled occurrence is discarded.  For
\(\pi:t\longrightarrow^*n\), define \(\cross^\downarrow(\pi,n)\) to be the
multiset of triples \((\ell,e,e')\) such that the contraction labelled
\(\ell\) stored event \(e\), and a residual of that node occurs in \(n\) with
final event field \(e'\).
\end{definition}

The labels are only a residual-accounting device; erasing them gives exactly
the reduction relation above.  The multiset records repeated contractions and
duplication, whereas \(\events(n)\) records set-valued final support.  Because
\(n\) is normal, every recorded residual event \(e'\) is normal.  For a
top-level closed record not moved through a binder, \(e'=\nf(e)\); the more
general residual formulation also covers contraction under binders followed
by substitution.

It would be incorrect to compare the raw event at contraction with final
support by syntactic equality.  For example, an event may be a beta-redex that
is retained by the record step and normalized only afterward.  The residual
trace records precisely the final representative that support can contain.

\section{Structural metatheory}\label{sec:metatheory}

\subsection{Effects as path semantics}

Every effect-bearing rule edge is labelled by the coefficient displayed in
its effect equation.  Ordinary computational premises have edge \(0\); the
domain edge of \(\textsc{Ty-Pi}\) and argument edge of \(\textsc{App}\) have
edge \(q\); and the stamp premise of a gate and the payload and event premises
of a node have edge \(1\).  Formation premises used only to validate a type
annotation are separate accepted judgments and contribute no term-effect
path.  Equivalently, their edge in the term derivation is \(\absent\).  An
\(\absent\)-edge kills a path.  For an atom occurrence path \(p\), let
\(\cost(p)\) be the \(\otimes\)-product, hence the maximum, of its edge labels.
A path is \emph{live} when its cost is below \(\absent\).

\begin{lemma}[Path characterization]\label{lem:path-characterization}
If a derivation has inferred effect \(\eps\), then for every atom \(z\),
\[
 \eps(z)=\bigoplus_{p:z\leadsto\mathrm{root}}\cost(p),
\]
where the empty sum is \(\absent\).
\end{lemma}

\begin{proof}
Induction on the last rule.  A variable or constant has the single empty path
of cost \(0\).  An ordinary constructor takes the minimum over the paths in
its premises.  Application prefixes every argument path by \(q\), producing
\(q\otimes\eta\).  Abstraction deletes paths beginning at its bound variable.
The gate and node cases prefix exactly the privileged paths by \(1\).  These
are precisely the effect equations in \cref{sec:type-rules,sec:term-rules}.
\end{proof}

\begin{theorem}[Principal and decidable effects]\label{thm:principal}
For a fixed well-formed signature, context, annotated expression, and result
role, its effect is unique when it is typable.  Effect inference and policy
validation are decidable.  For an unannotated lambda, elaboration uniquely
chooses the maximal admissible grade \(q=\theta(x)\), where \(\theta\) is the
principal body effect.  An explicit annotation is accepted exactly when
\(D(\sigma,\rho)\le q\le\theta(x)\).
\end{theorem}

\begin{proof}
The syntax determines the last rule.  Each effect equation uses finite
pointwise minimum, maximum, restriction, and unit-vector construction.
\textsc{Conv} does not change the effect, so the candidate principal effect is
computed independently of the later conversion decision.  The lambda equation
tests \(q\le\theta(x)\); maximal source elaboration chooses equality.  The
binder-policy test and ordinary policy validation are finite comparisons.
\end{proof}

\subsection{Substitution}

For \(\eps\) over \(\Gamma,x:_{\sigma}A\) and \(u\) with effect \(\eta\),
define
\begin{equation}\label{eq:single-subst}
  \eps\langle x:=\eta\rangle
  = (\eps\setminus x)\oplus(\eps(x)\otimes\eta).
\end{equation}
For a simultaneous substitution
\(\theta:\Delta\to\Gamma\), let \(M_\theta(y,z)\) be the effect at source
atom \(z\) of the term replacing \(y\).  Matrix action is over the min--max
semiring:
\begin{equation}\label{eq:matrix-subst}
 (M_\theta\star\eps)(z)
 =\bigoplus_{y\in\Gamma}\eps(y)\otimes M_\theta(y,z),
\end{equation}
with the signature-atom components of \(\eps\) included unchanged.

\begin{lemma}[Exact single substitution]\label{lem:substitution}
Suppose
\[
 \AFDeclSignature;\Gamma,x:_{\sigma}A\vdash_{\rho} J\triangleright\eps
 \quad\text{and}\quad
 \AFDeclSignature;\Gamma\vdash_{\sigma} u:A\triangleright\eta,
\]
where \(J\) is either a type or a term judgment.  Then
\[
 \AFDeclSignature;\Gamma\vdash_{\rho} J[u/x]
 \triangleright\eps\langle x:=\eta\rangle.
\]
\end{lemma}

\begin{proof}
Simultaneous induction on type formation and term typing.  The variable case
is the definition of \(\unit_x\).  Ordinary constructors use distributivity
of \(\otimes\) over \(\oplus\).  In the application case the induction
hypotheses and associativity of \(\otimes\) give
\[
 q\otimes(\eps(x)\otimes\eta)
 =(q\otimes\eps(x))\otimes\eta.
\]
Binder cases use alpha-renaming and effect restriction.  Gate and node cases
are the application calculation with fixed prefix \(1\).  In the lambda case
the bound name is chosen fresh for the substituend, so substitution does not
change its body-effect component and preserves \(q\le\theta(x)\).
Substitution in dependent result types uses the simultaneous type induction.
\end{proof}

\begin{theorem}[Policy-preserving substitution]\label{thm:substitution}
If both premises of \cref{lem:substitution} are policy valid and \(D\) is path
closed, then the substituted conclusion is policy valid.  Simultaneous
substitution has effect \(M_\theta\star\eps\), and composition of
substitutions is associative.
\end{theorem}

\begin{proof}
Consider a source role \(\tau\), the substituted variable role \(\sigma\),
and result role \(\rho\).  Validity gives effects at least
\(D(\tau,\sigma)\) and \(D(\sigma,\rho)\).  Their composite is at least
\[
 D(\tau,\sigma)\otimes D(\sigma,\rho)
 \ge D(\tau,\rho)
\]
by path closure.  Minimum over multiple substituted paths preserves the same
lower bound.  Associativity is min--max matrix associativity.
\end{proof}

Weakening is the special case that extends an effect by
\(\absent\).  Exchange is admissible for independent declarations, and
strengthening is admissible when the deleted component is \(\absent\).

\subsection{Preservation}

Effects are ordered pointwise.  Moving upward means that a dependency becomes
more mediated or disappears.

\begin{theorem}[Subject reduction with effect improvement]
\label{thm:preservation}
If
\(
 \AFDeclSignature;\Gamma\vdash_{\rho} t:A\triangleright\eps
\)
and \(t\longrightarrow t'\), then
\[
 \AFDeclSignature;\Gamma\vdash_{\rho} t':A\triangleright\eps'
 \qquad\text{for some }\eps\le\eps'.
\]
Consequently policy validity is preserved.  The analogous result holds for
reduction in types.
\end{theorem}

\begin{proof}
Compatible cases use monotonicity of \(\oplus\) and \(\otimes\).  For beta,
let \(\theta\) be the body effect.  The redex effect contains
\(q\otimes\eta\) with \(q\le\theta(x)\), while
\cref{lem:substitution} gives the contractum effect
\((\theta\setminus x)\oplus(\theta(x)\otimes\eta)\).  The two beta effects
are equal for maximally elaborated lambdas and otherwise the contractum effect
is greater.  Projection may discard one component and therefore only raises
its effect.  In the record step the origin-proof component disappears; all
other components retain exactly the same prefixes.

The only subtle compatible case is reduction beneath a lambda.  If its body
effect improves from \(\theta\) to \(\theta'\), then
\(q\le\theta(x)\le\theta'(x)\), so the same annotation remains admissible;
the external effect improves from \(\theta\setminus x\) to
\(\theta'\setminus x\).  This is the reason \textsc{Lam} uses a stable lower
bound rather than equality.  Type preservation follows from substitution in
the beta and projection cases and from the gate declaration's fixed result
type in the record case.
\end{proof}

\subsection{Normalization, confluence, and canonical forms}

The proofs in this subsection concern exactly the syntax in
\cref{sec:core}.  They do not cover arbitrary module rewrite rules.  We give
the reducibility argument because the record rule is not an instance of the
ordinary \(\AFDepFun\AFDepPair\) metatheory.

\paragraph{Simple skeletons.}
Erase indices, roles, and grades from a well-formed type but retain its outer
computational shape:
\[
 \tau ::= \iota_a\mid\iota_E\mid\iota_O
       \mid S\tau\mid\tau\Rightarrow\tau\mid\tau\times\tau .
\]
Write \(\lfloor A\rfloor\) for this skeleton.  Family applications map to
\(\iota_a\), \(\Event\) to \(\iota_E\), every origin type to
\(\iota_O\), stamps to \(S\lfloor A\rfloor\), dependent products to arrows,
and dependent sums to products.

\begin{lemma}[Skeleton invariance]\label{lem:skeleton}
If \(A\longrightarrow A'\), \(A\jdeq A'\), or \(A[u/x]\) is defined, then,
respectively,
\[
 \lfloor A\rfloor=\lfloor A'\rfloor,
 \qquad
 \lfloor A\rfloor=\lfloor A'\rfloor,
 \qquad
 \lfloor A[u/x]\rfloor=\lfloor A\rfloor.
\]
\end{lemma}

\begin{proof}
Reduction occurs only in recursive type positions and term indices; no rule
changes a type-former head.  The first and third statements follow by
structural induction.  The equality statement follows by induction over the
generation of definitional equality.
\end{proof}

Let \(\mathsf{SN}\) be the set of raw terms having no infinite reduction.
Define reducibility sets \(\mathcal R_\tau\) by induction on \(\tau\):
\begin{align*}
 \mathcal R_{\iota}&=\mathsf{SN},\\
 \mathcal R_{S\tau}&=
 \{t\in\mathsf{SN}\mid
   t\longrightarrow^*\mathsf{stamp}(u,e,o)
   \Rightarrow u\in\mathcal R_\tau\text{ and }e,o\in\mathsf{SN}\},\\
 \mathcal R_{\tau\Rightarrow\upsilon}&=
 \{t\in\mathsf{SN}\mid
   u\in\mathcal R_\tau\Rightarrow t\,u\in\mathcal R_\upsilon\},\\
 \mathcal R_{\tau\times\upsilon}&=
 \{t\in\mathsf{SN}\mid
   \fst(t)\in\mathcal R_\tau\text{ and }
   \snd(t)\in\mathcal R_\upsilon\}.
\end{align*}
A term is \emph{neutral} when it is not a lambda, pair, or stamp and has no
head redex.  Thus variables, public constants, private nodes, their neutral
spines, and a gate with a neutral credential are neutral; a gate whose
credential is visibly a stamp is not.

\begin{lemma}[Reducibility closure]\label{lem:candidate-closure}
For every skeleton \(\tau\):
\begin{enumerate}[label=(\roman*),leftmargin=2.5em]
\item \(\mathcal R_\tau\subseteq\mathsf{SN}\);
\item if \(t\in\mathcal R_\tau\) and \(t\longrightarrow t'\), then
\(t'\in\mathcal R_\tau\);
\item if \(r\) is neutral and every one-step reduct of \(r\) is in
\(\mathcal R_\tau\), then \(r\in\mathcal R_\tau\);
\item beta, projection, and record head expansion preserve reducibility when
their contracta and immediate computational components are reducible, and any
discarded proof component is strongly normalizing; and
\item stamps with reducible payload and strongly normalizing event and origin
fields are in \(\mathcal R_{S\tau}\).
\end{enumerate}
\end{lemma}

\begin{proof}
Simultaneous induction on \(\tau\).  Clause (i) is part of every defining
clause.  For (ii), an infinite reduction from \(t'\) would extend to one from
\(t\).  At an arrow, fix \(u\in\mathcal R_\sigma\); the step
\(t\,u\longrightarrow t'\,u\) and the induction hypothesis at the codomain
give the arrow clause.  At a product, use the corresponding steps under both
projections.  At a stamp skeleton, prepend \(t\longrightarrow t'\) to any
reduction of \(t'\) to a stamp.

For (iii), all immediate reducts are strongly normalizing, so the neutral
term is strongly normalizing.  At an arrow, fix
\(u\in\mathcal R_\sigma\).  The term \(r\,u\) is neutral.  Its reducts are
\(r'\,u\), handled by the premise and the arrow clause, or \(r\,u'\),
handled by subsidiary induction on the finite reduction height of \(u\).
The product case applies the same argument to \(\fst(r)\) and \(\snd(r)\).
For \(S\sigma\), any reduction from \(r\) to a stamp begins with a reduct
covered by the premise.  This proves neutral closure.

For (iv), induct on the result skeleton and, secondarily, on the sum of the
reduction heights of the immediate components.  A one-step reduct either is
the designated contractum or reduces a component and decreases that height;
clauses (ii) and (iii) close the cases.  For (v), every reduction of a stamp
is in one of its fields.  Their finite reduction heights give strong
normalization, and clause (ii) preserves the required component membership
at every canonical stamp reduct.
\end{proof}

\begin{lemma}[Gate and node reducibility]\label{lem:gate-reducibility}
Suppose all ordinary parameters of a registered gate are reducible at their
skeletons.
\begin{enumerate}[label=(\roman*),leftmargin=2.5em]
\item If \(t\), \(e\) are reducible at the payload and event skeletons, then
\(\mathsf{node}_g(\vec d;t,e)\in\mathcal R_{\lfloor B\rfloor}\).
\item If \(s\in\mathcal R_{S\lfloor A\rfloor}\), then
\(\mathsf{gate}_g(\vec d;s)\in\mathcal R_{\lfloor B\rfloor}\).
\end{enumerate}
\end{lemma}

\begin{proof}
For (i), use neutral closure and induction on the sum of the finite reduction
heights of \(\vec d,t,e\).  A node has no head step, and every reduct lowers
that sum.

For (ii), induct on the reduction heights of \(\vec d,s\).  Reducing a
parameter or the credential lowers the measure.  If the credential is
neutral, neutral closure applies.  If
\(s=\mathsf{stamp}(t,e,o)\), the definition of
\(\mathcal R_{S\lfloor A\rfloor}\) gives
\(t\in\mathcal R_{\lfloor A\rfloor}\) and \(e\in\mathsf{SN}\); the record
contractum is reducible by (i), and record head expansion from
\cref{lem:candidate-closure}(iv) gives the redex.  The origin field is strongly
normalizing and is discarded, so it creates no additional case.
\end{proof}

A substitution \(\gamma\) is \emph{reducible for \(\Gamma\)} when, in
context order, it maps each \(x:_{\rho}A\) to a term in
\(\mathcal R_{\lfloor A\gamma\rfloor}\).  Fresh variables give the identity
reducible substitution by neutral closure.

\begin{lemma}[Fundamental lemma]\label{lem:fundamental}
Let \(\gamma\) be reducible for \(\Gamma\).
\begin{enumerate}[label=(\roman*),leftmargin=2.5em]
\item If \(\AFDeclSignature;\Gamma\vdash_{\rho} A\;\jtype\), then \(A\gamma\) is
strongly normalizing.
\item If \(\AFDeclSignature;\Gamma\vdash_{\rho} t:A\), then
\(t\gamma\in\mathcal R_{\lfloor A\rfloor}\).
\end{enumerate}
\end{lemma}

\begin{proof}
Simultaneous induction on the accepted formation and typing derivations.
For formation, each type former has finitely many recursive type fields and
typed term indices.  The induction hypotheses make those fields strongly
normalizing; the multiset of their reduction heights strictly decreases at
every type step.  Under a binder, extend \(\gamma\) by the fresh neutral
variable.  This proves (i).

For (ii), variables are covered by reducibility of \(\gamma\), and constants
by neutral closure.  For a lambda, extend \(\gamma\) by an arbitrary
reducible argument and use the term induction hypothesis on the body; beta
head expansion in \cref{lem:candidate-closure} proves the arrow clause.
The grade inequality and role annotations are computationally inert.  For an
application use the arrow clause.  Pairing and projections use the product
clause and projection expansion.  A stamp uses
\cref{lem:candidate-closure}(v).  Gate and node terms use
\cref{lem:gate-reducibility}.  Conversion is harmless by
\cref{lem:skeleton}.  These are all term rules.
\end{proof}

\begin{theorem}[Strong normalization]\label{thm:sn}
Every well-typed \(\RATT\) type and term is strongly normalizing.
\end{theorem}

\begin{proof}
Use the identity reducible substitution in \cref{lem:fundamental}.
Term reducibility implies strong normalization by
\cref{lem:candidate-closure}(i), and the formation clause gives the type
case.
\end{proof}

\paragraph{Local peaks.}
We spell out the interaction between beta reduction and records rather than
appealing only to orthogonality.

\begin{lemma}[Residuals of one step]\label{lem:one-step-substitution}
If \(r\longrightarrow r'\), then
\(r[u/x]\longrightarrow^*r'[u/x]\).  If \(u\longrightarrow u'\), then
\(r[u/x]\longrightarrow^*r[u'/x]\), reducing once in every residual of
\(u\).  The same statements hold for substitution into types.
\end{lemma}

\begin{proof}
Structural induction on \(r\).  The variable case is immediate.  Binder
cases use alpha-renaming; every other case uses compatible closure.  The
second statement may require zero steps when \(x\) is absent and finitely many
steps when substitution duplicates \(u\).
\end{proof}

\begin{lemma}[Local confluence]\label{lem:local-confluence}
If a well-typed expression \(r\) reduces in one step to both \(r_1\) and
\(r_2\), then there is \(v\) with
\(r_1\longrightarrow^*v\longleftarrow^*r_2\).
\end{lemma}

\begin{proof}
Compare the two redex positions.  Disjoint positions commute.  Equal
positions use the same unique head rule.  It remains to consider a redex
strictly inside another head redex.

For beta, a body step is transported through substitution by the first part
of \cref{lem:one-step-substitution}; an argument step is transported to all
of its residuals by the second part.  A step in the lambda annotation is
discarded by both beta paths.  For a projection, a step in the selected
component is retained and commutes, while a step in the discarded component
is erased on both paths.  For a record, steps in gate parameters, payload,
or event are retained in the corresponding node fields and commute.  A step
in the origin proof is discarded by both paths.  No head rule overlaps a
different head rule: lambda, pair, and gate-with-stamp have distinct heads.
Compatible reductions in type indices and annotations are the same retained
or discarded context cases.  This exhausts the grammar.
\end{proof}

\begin{theorem}[Confluence]\label{thm:confluence}
Typed reduction in \(\RATT\) is confluent.
\end{theorem}

\begin{proof}
By \cref{thm:sn} every typed reduction relation is terminating, and by
\cref{lem:local-confluence} it is locally confluent.  Newman's lemma therefore
gives confluence.
\end{proof}

\paragraph{Canonical forms.}
For normal annotations and fields, define neutral and normal terms mutually:
\begin{align*}
 h ::={}&x\mid c\mid
   \mathsf{node}_g(\vec n;n,e)\mid
   \mathsf{gate}_g(\vec n;r),\\
 r ::={}&h\mid r\,n\mid\fst(r)\mid\snd(r),\\
 n ::={}&r\mid\lambda^q(x:_{\sigma}A).n
       \mid(n,n)\mid\mathsf{stamp}(n,e,o).
\end{align*}
In the gate head, the credential \(r\) is neutral, so the record redex is
absent.  Event and origin fields in the grammar are themselves normal terms.

\begin{theorem}[Canonical forms]\label{thm:canonical-forms}
Every well-typed normal term has exactly one of the preceding normal forms.
Moreover, for closed normal terms:
\begin{enumerate}[label=(\roman*),leftmargin=2.5em]
\item a term of \(\Stamped_\sigma(A)\) is a stamp;
\item a term of \(\Event\) is an event constant from a unique certificate or
attestation entry, and a term of \(\Origin_A(e,t)\) is that entry's matching
origin constant;
\item a term of \(\AFDepFun^q(x:_{\sigma}A).B\) with
\(q\in\{1,\absent\}\) is a lambda; and
\item a closed neutral at a protected grade-free type is headed by an
ordinary public constant or a private node; it is never headed by a public
gate.
\end{enumerate}
\end{theorem}

\begin{proof}
The grammar statement is structural induction on the term.  A normal
application cannot have a lambda head, a normal projection cannot have a pair
head, and a normal gate cannot have a stamp credential; these are exactly the
four omitted redex shapes.

For (i), a closed neutral cannot be variable-headed.  Primitive and
certificate-payload declarations have ordinary ultimate result heads,
attestations introduce no payload constant, gates have protected results, and
nodes have their gates' protected results.  Hence none has stamped type.  The
only remaining normal introduction is a stamp.
For (ii), the signature rules prohibit primitive, gate, or payload constants
with ultimate head \(\Event\) or \(\Origin\); certificate and attestation
entries are their only sources.  Fresh qualified names give unique ownership.  Type-head
injectivity, which follows from \cref{lem:skeleton} and confluence, forces an
origin constant's event and payload to match the expected normal fields.

For (iii), a closed neutral head of positive or absent function grade cannot
be a primitive or certificate payload because their types are grade-free;
an attestation adds only event and origin constants, and a gate or node has a
grade-free result.  Thus the normal form is a lambda.  For (iv), a closed gate credential is a stamp by (i), making the
gate a redex.  The remaining closed neutral heads are precisely the two
listed classes; grade-freedom excludes a nonzero latent application edge in
their spine.
\end{proof}

\begin{corollary}[Canonical credentials]\label{cor:credentials}
Every closed normal term of a stamped type is headed by
\(\mathsf{stamp}\).  Its event and origin come from one unique evidence entry.
For a certificate its payload is the fresh payload constant; for an exact
attestation its payload is definitionally equal to the old closed term named
by that entry.
\end{corollary}

\begin{proof}
Apply \cref{thm:canonical-forms}(i--ii) and type-head injectivity for
\(\Origin_A(e,t)\).  Its normalized indices distinguish the certificate and
exact-attestation cases.
\end{proof}

\begin{theorem}[Decidable checking and conversion]\label{thm:decidability}
For a finite well-formed signature, context formation, type checking, effect
inference, policy validation, and definitional equality are decidable.
\end{theorem}

\begin{proof}
Use the bidirectional organization specified in \cref{app:checker}: synthesize
neutral forms and check introductions against an expected type.  All role,
grade, freshness, effect, and policy conditions range over finite data; the
lambda test is the decidable comparison \(q\le\theta(x)\).  By
\cref{thm:sn,thm:confluence}, two well-typed expressions are definitionally
equal exactly when their normal forms are alpha-equivalent.  Normalization and
the structural comparison therefore terminate.
\end{proof}

\section{Gate adequacy}\label{sec:adequacy}

\subsection{Static path mediation}

\begin{theorem}[Path mediation]\label{thm:mediation}
Let a policy-valid type or term derivation have result role \(\rho\) and
principal effect \(\eps\).  For every source atom \(z\) of role \(\sigma\):
\begin{enumerate}[label=(\roman*),leftmargin=2.5em]
\item if \(D(\sigma,\rho)=1\), every live effect path from \(z\) to the root
contains a registered gate or node edge, a grade-\(1\)
\textnormal{\textsc{Ty-Pi}} domain
edge, or a latent grade-\(1\) \textnormal{\textsc{App}} edge; and
\item if \(D(\sigma,\rho)=\absent\), there is no live effect path from \(z\)
to the root: every such derivation path contains an \(\absent\)-edge.
\end{enumerate}
These path statements are stable under policy-valid substitution.  Reduction
may discard paths but cannot create a path below the policy demand.
\end{theorem}

\begin{proof}
By \cref{lem:path-characterization}, \(\eps(z)\) is the minimum path cost.
Validity with demand \(1\) implies that every live cost is at least \(1\).
Inspection of the effect equations shows that only a gate edge, node edge,
grade-\(1\) \textsc{Ty-Pi} edge, or grade-\(1\) \textsc{App} edge can raise a
live path from \(0\) to \(1\).  Demand \(\absent\) implies that every path has
cost \(\absent\), hence contains an \(\absent\)-edge.  Stability is
\cref{thm:substitution,thm:preservation}.
\end{proof}

The second clause is an effect statement, not syntactic erasure.  For example,
if \(\vdash_{\Ren}N\;\jtype\) and \(\vdash_{\Frm}B\;\jtype\), then the normal type
\(
  \AFDepFun^{\absent}(x:_{\Ren}N).B
\)
still contains \(N\), although its occurrence has no live path.  Likewise, a
bare \(\AFDepFun^1\)-type supplies a type-formation edge but has no term body to
summarize.  Operational realization therefore requires a term-level
canonical-form statement.

\begin{lemma}[Positive-function realization]\label{lem:positive-function}
Every closed normal term of type
\(
  \AFDepFun^1(x:_{\sigma}A).B
\)
at result role \(\rho\) is a lambda \(\lambda^1(x:_{\sigma}A).t\), and the principal effect of \(x\) in
\(t\) is either \(1\) or \(\absent\).  If \(x\) occurs on a live body path,
that effect is \(1\).  Consequently, applying such a closed function exposes
either the gated body path summarized by its grade or an argument that the
body discards.
\end{lemma}

\begin{proof}
A closed normal function cannot be headed by a variable.  Primitive and
certificate declarations are grade-free and have ordinary ultimate heads;
gate and node results are grade-free.  Hence none can have the displayed
positive function type.  The only remaining canonical form is a lambda, and
\textsc{Lam} requires \(1\le\theta(x)\).  The only possibilities are \(1\)
and \(\absent\), and liveness excludes the latter.
\end{proof}

The word ``closed'' in \cref{lem:positive-function} does not eliminate neutral
applications beneath an unapplied outer lambda.  For example, a bound variable
may itself be assumed to have positive function type.  The dynamic theorem
below therefore concerns grade-free result types, where such latent positive
interfaces cannot remain in the observed normal form.

Call a public declaration occurrence in an expression \emph{effect relevant}
when its path in the principal derivation is live in the sense of
\cref{sec:metatheory}, that is, when its path cost is below \(\absent\).
Bound variables are handled by the binder-policy premise but are not public
extension atoms.  Occurrences confined to type annotations behind
\(\absent\)-edges are syntactically present but are not effect relevant.

\begin{lemma}[No latent protected escape]\label{lem:no-latent-escape}
Let \(\vdash_{\rho}n:A\) be closed and normal, where
\(\rho\in\{\Frm,\Prs\}\) and \(\mathsf{gradefree}(A)\).  Outside private
nodes:
\begin{enumerate}[label=(\roman*),leftmargin=2.5em]
\item no effect-relevant public declaration path can end in a grade-
\(1\) or grade-\(\absent\) neutral application edge; and
\item a normal public gate cannot carry an effect-relevant closed public
declaration occurrence in its credential.
\end{enumerate}
\end{lemma}

\begin{proof}
Use the normal/neutral grammar of \cref{thm:canonical-forms}, inducting
simultaneously on \(n\) and the grade-free type through which it is checked.
For (i), an application with nonzero latent edge cannot have a lambda head,
since it would be a beta-redex.  Primitive and certificate types are
grade-free, and gate and node result types are grade-free.  Its neutral head
must therefore be a bound variable whose declared type contains that nonzero
grade.  The type is a subexpression of an enclosing binder domain, contradicting
the recursive definition of \(\mathsf{gradefree}(A)\).

For (ii), a normal gate has a neutral, rather than stamped, credential.
No public declaration may have ultimate stamped result, and no gate or node
has witness result.  Its neutral head is therefore a bound variable whose
ultimate result role is \(\Wit\).  Along the enclosing protected lambda,
\textsc{Ty-Pi} requires
\(D(\Wit,\rho)\le q\), hence \(1\le q\), contradicting grade-freedom.
The same argument applies when the public occurrence is in an ordinary
argument of that neutral credential head.  The induction passes through
ordinary grade-zero lambdas, pairs, and neutral spines; those constructors
introduce no nonzero edge.
\end{proof}

\begin{lemma}[Grade-free normal-form mediation]\label{lem:gradefree-mediation}
Let \(\vdash_{\rho}n:A\) be closed and normal, with
\(\rho\in\{\Frm,\Prs\}\) and \(\mathsf{gradefree}(A)\).  Every
effect-relevant occurrence in \(n\) of a public atom \(z\) of role \(\sigma\) satisfying
\(D(\sigma,\rho)=1\) lies below a private node edge.  No occurrence with
policy demand \(\absent\) is effect relevant in \(n\).
\end{lemma}

\begin{proof}
By \cref{thm:mediation}, a live demand-\(1\) occurrence path contains a node
edge, a public-gate edge, or a nonzero latent application edge; formation
premises for type annotations are not term-effect paths.  The latter two
possibilities outside nodes are excluded by
\cref{lem:no-latent-escape}.  Induction through the canonical grammar then
places the occurrence below the first private node on its path.  For demand
\(\absent\), path mediation says every derivation path has an
\(\absent\)-edge and hence is not live.
\end{proof}

The restrictions on primitive grades and gate result types are essential.  If
a module could simply declare \(c:_{\Frm}\AFDepFun^1(x:_{\Wit}A).B\), the type would
assert mediation without any term that emits a record.  Likewise, if a gate
node could have a result containing an opaque positive-grade function, its
existing record could be misused to justify unrelated future arguments.
\(\RATT\) therefore admits a closed canonical function of positive grade only
as a lambda.  Maximal source elaboration gives that lambda its exact principal
parameter grade; after reduction, a stable positive annotation may conservatively
remain on a body that has become unused, but every live parameter path still
has grade \(1\).  Open terms may assume positive function variables, which is
why the dynamic result has a grade-free observation boundary.  A registered
gate remains a saturated term former that such a body may invoke.

\subsection{Dynamic records}

\begin{lemma}[Record persistence]\label{lem:record}
If \((\ell,e,e')\in\cross^\downarrow(\pi,n)\), then a residual node labelled
\(\ell\) occurs in \(n\), its final event field is \(e'\), and
\(e'\in\events(n)\).  The field \(e'\) is obtained from \(e\) by the reductions
and substitutions along that residual; if \(e\) is closed at contraction, then
\(e'=\nf(e)\).  Before normalization, every reduct after the labelled
contraction contains at least one residual of that node unless a surrounding
redex discards every residual of the protected subterm.
\end{lemma}

\begin{proof}
The record rule inserts \(e\) in the labelled node.  No rule removes an event
field from a surviving node.  Compatible reduction may reduce or substitute
inside that field, move the node, or duplicate it through substitution.  It
can remove all residuals only by discarding the containing subterm.  The
definition of \(\cross^\downarrow\) records the final event field \(e'\) of
each surviving residual, and the structural definition of \(\events\) includes
that field.  Residual propagation applies exactly the ambient reductions and
substitutions to the stored field.  When \(e\) is closed, substitution has no
effect; normality of \(n\) and confluence give \(e'=\nf(e)\).
\end{proof}

\begin{lemma}[Public node ancestry]\label{lem:node-ancestry}
If \(t\) is public and \(\pi:t\longrightarrow^*n\), every private node
occurrence in \(n\) is a residual of at least one record contraction in
\(\pi\).
\end{lemma}

\begin{proof}
Induct over \(\pi\).  The initial expression has no node.  Ordinary beta and
compatible steps can move, duplicate, reduce inside, or discard existing node
residuals but cannot create a new node head.  Projection can only retain or
discard them.  The record rule is the unique rule whose contractum introduces
a node, and its fresh label identifies every later copy.  Hence every final
node has the asserted ancestor.
\end{proof}

\begin{theorem}[Static--dynamic gate adequacy]\label{thm:adequacy}
Let \(\vdash_{\rho}t:A\) be a closed public, policy-valid term, with
\(\rho\in\{\Frm,\Prs\}\) and \(\mathsf{gradefree}(A)\), and let
\(\pi:t\longrightarrow^*n\) be a normalization sequence.  Then:
\begin{enumerate}[label=(\roman*),leftmargin=2.5em]
\item every effect-relevant occurrence in \(n\) whose policy demand is \(1\)
lies beneath a residual node created by a registered gate contraction in
\(\pi\); and
\item if \((\ell,e,e')\in\cross^\downarrow(\pi,n)\), then
\(e'\in\events(n)\).
\end{enumerate}
In particular, a successful exceptional crossing cannot be both
computationally relevant to the result and absent from its normal provenance.
\end{theorem}

\begin{proof}
By \cref{lem:gradefree-mediation}, every relevant positive occurrence in
\(n\) lies below a private node.  Since \(t\) is public, every such node has at
least one record-step ancestor in \(\pi\) by
\cref{lem:node-ancestry}; residual labels make this ancestry explicit even
when substitution duplicates the node.  This proves (i).
Part (ii) is \cref{lem:record}.  Confluence makes \(n\), and hence its
structural event support, independent of the chosen normalization sequence,
up to alpha-equivalence.
\end{proof}

The qualification ``contributing to the result'' is necessary.  A gate may be
constructed in an argument that an enclosing function ignores.  Normalization
correctly removes both that computation and its record.

\section{Modules and extension theorems}\label{sec:modules}

\subsection{Module forms}

A module over \(\AFDeclSignature\) is a finite well-formed signature suffix
\(\mathcal M\) whose public names share an owner and are fresh for \(\AFDeclSignature\).
Concatenation is written \(\AFDeclSignature+\mathcal M\).  The core distinguishes the
following syntactic classes; membership is decidable by inspection of the
suffix.

\begin{description}[leftmargin=3.5cm,style=nextline]
\item[Rendering-only]
Every new family and primitive has result role \(\Ren\).  There is no
certificate bundle, exact attestation, or gate.
\item[Fresh-witness-family]
Every new family has role \(\Wit\), and every new primitive is a constructor
whose result is headed by one of those new families.  There is no certificate
bundle, exact attestation, or gate.
\item[Evidence-inhabitant]
The nonempty suffix consists of any mixture of certificate bundles for old
closed types at role \(\Wit\) and exact attestations of old closed public
terms at old closed witness-data types.  It adds no family, ordinary
primitive, or gate.  An attested payload term is old.
\item[Gate module]
The suffix contains one or more registered gates and may therefore add new
protected construction power.  Its guarantees are mediation, record
adequacy, normalization, and decidable checking, not conservativity.
\end{description}

A script generator is an ordinary primitive at role \(\Scr\).  Script
equations are not a declaration form in \(\RATT\); they belong to a separate
typed setoid or categorical presentation.  This prevents an extension from
collapsing old script equality through fresh intermediaries.

\subsection{Old terms and equality reflection}

An expression is \emph{old} when all of its public declaration atoms belong to
\(\AFDeclSignature\).  An extension is term-conservative at an old closed type \(A\)
of role \(\rho\) when every closed public extended term of that type is
definitionally equal to an old term.  It reflects equality when equality of
two old public terms in the extension implies their equality in the base.
Public reachability is essential: the internal \textsc{Node} rule is a
reduction target, not a source-language constructor.

The absence of internal universes is used exactly here.  A new family symbol
cannot be packaged as data in an old type, and an old gate payload cannot be
instantiated by a family declared only later.

\begin{lemma}[Old-head and evidence-entry inversion]\label{lem:extension-inversion}
Let \(\AFDeclSignature+\mathcal M\) add no protected family, protected primitive, or
gate.
\begin{enumerate}[label=(\roman*),leftmargin=2.5em]
\item A new rendering primitive has ultimate head a new rendering family; a
new fresh-family constructor has ultimate head its new witness family.  It
cannot head a normal term at an old type.
\item If \(\mathcal M\) adds neither a certificate bundle nor an exact
attestation, every closed normal stamp at an old stamped type is old.
\item If \(\mathcal E\) is evidence-inhabitant and a closed normal stamp at
an old stamped type contains a new atom, then for one unique evidence entry
its normal fields have one of the forms
\[
 \mathsf{stamp}(a_d,e_d,o_d)
 \quad\text{or}\quad
 \mathsf{stamp}(\nf(t_q),e_q,o_q),
\]
for a certificate or exact attestation, respectively.
\item If a public closed term reduces to a node residual containing an
\(\mathcal E\)-atom in a retained payload or event, its normal event field is
owned by \(\mathcal E\).  A new retained payload atom occurs only in the
certificate case and is paired with that certificate's event.
\end{enumerate}
\end{lemma}

\begin{proof}
For (i), signature order prevents an old declaration type from mentioning a
new family.  Type-head injectivity follows because reduction never changes a
type-former head and confluence gives unique normal types.  Hence a neutral
spine headed by either new declaration cannot synthesize an old result head.

For (ii), \cref{thm:canonical-forms} makes the credential a stamp and its
closed normal origin field an evidence-entry origin constant.  Because the
extension adds no certificate or attestation entry, that origin is old.
Injectivity of its normal \(\Origin\) type forces the event and payload fields
to be old as well.

For (iii), the same canonical-form argument applies.  The evidence extension
has no new constants other than certificate triples and attestation
event--origin pairs.  If the origin constant is \(o_d\), its declared type is
\(\Origin_A(e_d,a_d)\), so injectivity forces the stamp's event and payload
to \(e_d\) and \(a_d\).  If it is \(o_q\), its type is
\(\Origin_A(e_q,t_q)\), so the normal event and payload are \(e_q\) and
\(\nf(t_q)\).  The module premise makes \(t_q\) old, and qualified freshness
makes the entry unique.

For (iv), use \cref{lem:node-ancestry} to trace the residual to its record
step.  Compose the substitutions encountered by that residual between the
contraction and the closed final term, and apply the composite substitution to
the creating stamp, including its then-discarded origin field.  This produces
a closed well-typed stamp whose retained payload and event normalize to the
final residual fields.  Confluence and (iii) align a new payload atom with its
certificate event.  In the attestation case the payload is old and the new
retained atom is the event itself.  Substitution may duplicate the residual
and reduction may normalize its fields, but neither changes the qualified
event constant; \cref{lem:record} carries it to the final node.
\end{proof}

\begin{lemma}[Grade-free extension support]\label{lem:extension-support}
Let \(A\) at role \(\rho\) be old, closed, and grade-free, let \(\vdash_{\rho}t:A\) be closed
and public in an extension that adds no protected family, protected primitive,
or gate, and put \(n=\nf(t)\).
\begin{enumerate}[label=(\roman*),leftmargin=2.5em]
\item If the extension is rendering-only and \(\rho\ne\Ren\), then \(n\) is
old.
\item If it is fresh-witness-family and \(\rho\in\{\Frm,\Prs\}\), then \(n\)
is old.
\item If it is evidence-inhabitant and
\(\rho\in\{\Frm,\Prs\}\), every new atom in \(n\) occurs in a residual of an
old gate node, and that node's event field is a qualified event constant owned
by the extension.
\end{enumerate}
\end{lemma}

\begin{proof}
Induct on the canonical grammar of \(n\), simultaneously on the old
grade-free type against which each subterm is checked.  By
\cref{lem:extension-inversion}(i), a new rendering or fresh-witness
declaration cannot be the head at an old type.  A new atom could therefore
survive only as an argument of an old neutral head, in a stamp, or in a private
node.  Grade-free latent-edge exclusion
(\cref{lem:no-latent-escape}) removes the first possibility whenever the
policy requires a nonzero edge.  Ordinary grade-zero lambdas and pairs pass
the induction hypothesis to their old component types.

For (i), a rendering occurrence has demand \(\absent\) into every
non-rendering result.  A closed old gate credential is old by
\cref{lem:extension-inversion}(ii), and
\cref{lem:gradefree-mediation} says it has no live path.  The only syntactic
place that could retain it would be an \(\absent\)-graded neutral application,
which \cref{lem:no-latent-escape} excludes.  Thus \(n\) is old.

For (ii), a fresh witness occurrence has demand \(1\) into a protected
result, so \cref{lem:gradefree-mediation} places it below an old node.  Trace
that node to its creating old gate using \cref{lem:node-ancestry}.  The
creating credential is at an old stamped type and is old by
\cref{lem:extension-inversion}(ii).  Node parameters are handled by the
outer induction, so no fresh occurrence can enter a retained field.  Therefore
\(n\) is old.

For (iii), mediation places every new payload, event, or origin occurrence
beneath an old node.  Origin fields are discarded by the record step, and
\cref{lem:extension-inversion}(iv) says that the nearest residual retaining a
new payload has the matching extension-owned event.  New event occurrences
are already such retained event fields.  Hence every new atom has a witnessing
old node with an \(\mathcal E\)-owned event field.
\end{proof}

\begin{remark}[Why grade-freedom is necessary]\label{rem:gradefree-boundary}
The restriction is not cosmetic.  Let \(\vdash_{\Ren}G\;\jtype\),
\(\vdash_{\Wit}W\;\jtype\), and \(\vdash_{\Frm}B\;\jtype\) be old.  Put
\[
 F_R=\AFDepFun^{\absent}(x:_{\Ren}G).B,
 \qquad F_W=\AFDepFun^1(x:_{\Wit}W).B.
\]
A rendering-only extension adding \(r:G\) creates the closed normal, non-old
term
\(
 \lambda^0(f:_{\Frm}F_R).f\,r
\)
of the old non-rendering type \(\AFDepFun^0(f:_{\Frm}F_R).B\).  Likewise, an evidence
bundle adding \(a_d:W\) creates
\(
 \lambda^0(f:_{\Frm}F_W).f\,a_d
\)
at an old form-role type without yet producing a record.  Both dependencies
are latent behind a neutral higher-order call, and both result types contain a
nonzero grade.  Thus unrestricted higher-order conservativity and traceability
would be false.
\end{remark}

\begin{theorem}[Rendering noninterference]\label{thm:rendering}
Let \(\mathcal R\) be rendering-only.  For every old closed type
\(A\) of role \(\rho\), with \(\rho\ne\Ren\) and \(\mathsf{gradefree}(A)\), the extension
\(\AFDeclSignature+\mathcal R\) is term-conservative at \(A\) and reflects equality of
old terms.  In particular, it cannot change inhabitation of an old form,
presentation, witness or script type at this observation boundary.  This is
the book's key idea, that nothing may depend on a reading, as a theorem about
extensions.
\end{theorem}

\begin{proof}
Normalize a closed public extended \(\vdash_{\rho}t:A\).  By
\cref{lem:extension-support}(i), \(\nf(t)\) is old, proving term
conservativity.

An old term contains no new atom and hence has exactly its base reductions.
If two old terms are equal in the extension, confluence gives a common normal
form, which is reached using only base reductions.  Their equality is
therefore reflected.
\end{proof}

\begin{theorem}[Fresh-family conservativity]\label{thm:fresh}
Let \(\mathcal F\) be a fresh-witness-family module.  It is term-conservative
and equality-reflecting at every old closed grade-free type whose role is
\(\Frm\) or \(\Prs\).
\end{theorem}

\begin{proof}
Term conservativity is \cref{lem:extension-support}(ii).  Equality reflection
is the same old-reduction argument as in \cref{thm:rendering}.
\end{proof}

The protected-role qualification is exact.  A fresh witness family is new
vocabulary at role \(\Wit\), and the policy freely permits witnesses to occur
in scripts and renderings.  No conservativity claim is made there.

\subsection{Evidence inhabitants}

For a module \(\mathcal E\), write \(\Event_{\mathcal E}\) for its qualified
event constants.  Each \(e\in\Event_{\mathcal E}\) belongs to one unique
evidence entry and has one matching indexed origin owned by \(\mathcal E\).
A certificate pairs it with a fresh payload constant; an exact attestation
pairs it with the old closed term named by the entry.

\begin{theorem}[No silent enabling]\label{thm:no-silent}
Let \(\mathcal E\) be an evidence-inhabitant module and let
\(\vdash_{\rho}t:A\) be closed and public in \(\AFDeclSignature+\mathcal E\), where \(A\) is
old and grade-free and \(\rho\in\{\Frm,\Prs\}\).  If \(\nf(t)\) is not old,
then
\[
  \events(\nf(t))\cap\Event_{\mathcal E}\ne\varnothing.
\]
The extension still reflects judgmental equality of old terms.
\end{theorem}

\begin{proof}
By \cref{lem:extension-support}(iii), any new atom in \(\nf(t)\) occurs below
an old node whose normal event field is owned by \(\mathcal E\).  The
definition of \(\events\) therefore gives the displayed nonempty intersection.
The module adds no reduction rule, so old equality is reflected as before.
\end{proof}

This theorem replaces the false claim that witness-only extension preserves
old forms.  It states the stronger property appropriate to conditional
construction: the extension may enable an old gate, but the enabled result
identifies the responsible module in normal form.

\subsection{Composition}

\begin{theorem}[Module composition]\label{thm:composition}
If \(\mathcal M\) is well formed over \(\AFDeclSignature\) and \(\mathcal N\) is well
formed over \(\AFDeclSignature+\mathcal M\), with disjoint qualified names, then
\(\mathcal M;\mathcal N\) is well formed over \(\AFDeclSignature\).  Concatenation is
associative.  Principal effects, normalization, mediation, and record
adequacy are preserved.  A conservativity corollary is exported only when the
combined suffix remains in the corresponding rendering-only or fresh-family
class, and only at the grade-free observation boundary of
\cref{thm:rendering,thm:fresh}.
\end{theorem}

\begin{proof}
Signature formation is an ordered telescope, so the first claim and
associativity are immediate.  The effect algebra and policy are global and do
not change under extension.  Qualified freshness keeps evidence-entry
ownership disjoint.  The metatheorems depend only on well-formedness.  Syntactic module
classes are not closed under arbitrary concatenation, which is why their
conservativity conclusions are recomputed rather than assumed.
\end{proof}

\section{Erasure and provenance}\label{sec:erasure}

\subsection{Role erasure}

Let \(K_\AFDeclSignature\) be the role-free dependent \(\AFDepFun\AFDepPair\)-calculus with a
unit type, obtained from the same public families, ordinary primitives, and
certificate payload constants.  Exact attestations add no payload
declaration.  Event and origin constants from both evidence-entry forms map to unit.
Each gate contributes
an ordinary erased constructor
\[
  \overline g:|\Delta|\to|A|\to|B|.
\]
Erasure removes role and grade annotations and is homomorphic on ordinary
syntax.  On provenance syntax it is defined by
\[
\begin{array}{rclcrcl}
|\Event|&=&\mathbf 1,
&\qquad&|\Origin_A(e,t)|&=&\mathbf 1,\\
|\Stamped_\sigma(A)|&=&|A|,
&&|\mathsf{stamp}(t,e,o)|&=&|t|,\\
|\mathsf{gate}_g(\vec d;s)|&=&
 \overline g(|\vec d|;|s|),
&&|\mathsf{node}_g(\vec d;t,e)|&=&
 \overline g(|\vec d|;|t|).
\end{array}
\]
Event and origin constants erase to the sole unit term.  The target contains
the ordinary beta and projection rules and no record rule.

\begin{theorem}[Erasure simulation]\label{thm:erasure}
If \(\AFDeclSignature;\Gamma\vdash_{\rho} t:A\), then
\(|\Gamma|\vdash|t|:|A|\) in \(K_\AFDeclSignature\).  If
\(t\longrightarrow t'\), then either
\(|t|\longrightarrow|t'|\) or \(|t|=|t'|\).  Hence definitional equality is
preserved by erasure.
\end{theorem}

\begin{proof}
Induction on typing and reduction.  Ordinary cases are homomorphic.  Stamp
typing erases to payload typing.  The gate and node rules both erase to the
typing rule of \(\overline g\).  Beta and projections reduce in the target;
the record step maps both sides to the same erased constructor application.
\end{proof}

Erasure is intentionally non-faithful.  Two normal nodes with different event
constants may have identical erasures.  This loss is not an unspecified
semantic accident: provenance loss occurs at the displayed event, origin,
stamp, gate, and node clauses.  Erasure may additionally identify expressions
that differ only in role or grade annotations; no claim of a complete kernel
characterization is made here.

\subsection{Provenance normal forms}

For a closed protected term \(t\), define
\[
  \mathsf{pnf}(t)=\nf(t),
  \qquad
  \mathsf{support}(t)=\events(\nf(t)).
\]
Because private nodes are intrinsically typed by their gate result, a
provenance normal form cannot contain a record at an unrelated result type.

\begin{proposition}[Provenance stability]\label{prop:prov}
If closed protected terms \(t\) and \(u\) are definitionally equal, then
\[
 \mathsf{pnf}(t)=\mathsf{pnf}(u)
 \quad\text{and}\quad
 \mathsf{support}(t)=\mathsf{support}(u).
\]
Their erasures may agree even when their provenance normal forms differ.
\end{proposition}

\begin{proof}
Confluence gives the same normal form for definitionally equal terms.  Event
support is a deterministic structural function on that normal form.  The last
claim follows from the event-forgetting erasure clauses.
\end{proof}

\section{The book in the calculus}\label{sec:artforms}

The core calculus is independent of a domain signature.  This section
declares the book's frames in it: forms with their construction history,
installed presentations, interpretation witnesses, states with their
legality, scripts with their action, and renderings \cite{Book}.  The
declarations are an instantiation schema rather than a formalization of the
whole book.  The octagonal-rosette module of \cref{sec:rosette-module} and the
interpretation-profile instance of \cref{sec:interpretation-witnesses} are
finite admitted modules whose derivations, reduction traces, extensions and
erasure are worked out against the printed rules; \cref{sec:capture} writes
the capture of \cref{sec:chess} in the calculus.

\subsection{Forms and construction history}

Declare at role \(\Frm\)
\[
 \vdash_{\Frm}\AFFinCode\;\jtype,
 \qquad
 A:_{\Frm}\AFFinCode\ \vdash_{\Frm}\ \Hist(A)\;\jtype
\]
with ordinary constructors for empty, singleton, finite-index, and dependent
assembly histories.  Define
\[
  \vdash_{\Frm}\Form\;\jtype,
  \qquad
  \Form\coloneqq\AFDepPair(A:_{\Frm}\AFFinCode).\Hist(A).
\]
A form retains its carrier and its construction history, as the book's
\FORM{} does.  All generic form constructors use grade \(0\) arguments and
therefore add no special authority.

In the context \(X:_{\Frm}\Form\), let
\(X:_{\Frm}\Form\vdash_{\Wit}\mathsf{AttachData}_k(X)\;\jtype\) contain an
attaching shape, map, and boundary evidence.  Attachment is a registered
form gate
\[
 \mathsf{attach}_k:
 (X:_{\Frm}\Form);\,
 \Stamped_\Wit(\mathsf{AttachData}_k(X))
 \Longrightarrow_{\Frm}\Form.
\]
Its private normal constructor has the schematic form
\[
 \mathsf{node}_{\mathsf{attach}_k}(X;a,e),
\]
where \(a\) and \(e\) have grade-\(1\) edges.  Carrier evaluation may erase to
generic assembly and a finite quotient, while the form's normal form retains
the attachment event.  Quotient normalization itself is not a core RATT rule
and requires a separate certified computational extension.

For a closed \(X:\Form\) and an admitted certificate bundle
\[
 d=(a_d,e_d,o_d):\mathsf{AttachData}_k(X),
\]
the rules derive the public term and its record step
\[
 \vdash_{\Frm}t:\Form,
 \qquad
 t=\mathsf{gate}_{\mathsf{attach}_k}
   (X;\mathsf{stamp}(a_d,e_d,o_d))
 \longrightarrow
 \mathsf{node}_{\mathsf{attach}_k}(X;a_d,e_d).
\]
Before contraction its principal effect is
\[
 \unit_{\mathsf{attach}_k}\oplus\unit_X\oplus
  1\otimes(\unit_{a_d}\oplus\unit_{e_d}\oplus\unit_{o_d});
\]
after contraction the \(\unit_{o_d}\) component is discarded and the other
components retain their prefixes.  The reduct is normal when its fields are
normal, and \(e_d\in\events(\nf(t))\).  This is the concrete Artforms instance
of subject reduction with effect improvement and residual record adequacy.

\subsection{A worked octagonal-rosette module}\label{sec:rosette-module}

We now give one finite module rather than only a declaration schema.  It uses
an evidence extension to authorize an octagonal rosette inlay and a
rendering-only extension to observe the result.  All displayed declarations
are instances of the rules in \cref{app:rules}.

Start from the preceding Artforms signature and add ordinary constants
\[
 c_8:_{\Frm}\AFFinCode,
 \qquad
 h_8:_{\Frm}\Hist(c_8),
 \qquad
 X_8:_{\Frm}\Form,
 \qquad X_8\coloneqq(c_8,h_8).
\]
The constants have family-headed result types, while \(X_8\) is a derived
pair.  Add the witness family
\[
 \mathsf{RosetteSpec}_8:
 (X:_{\Frm}\Form)\Rightarrow_{\Wit}\Type,
\]
whose intended fields describe eight boundary anchors, cyclic order, and
nonintersection evidence.  The registered gate is
\[
 \mathsf{inlay}_8:
 (X:_{\Frm}\Form);\,
 \Stamped_\Wit(\mathsf{RosetteSpec}_8(X))
 \Longrightarrow_{\Frm}\Form
 \quad[\mathsf{node}_{\mathsf{inlay}_8}].
\]
Its ordinary parameter has demand \(0\), its credential edge has grade \(1\),
and its result \(\Form\) is grade-free.

The evidence-inhabitant extension \(\mathcal E_8\) is the single certificate
bundle
\[
 d_8=(a_8,e_8,o_8):\mathsf{RosetteSpec}_8(X_8).
\]
Before \(\mathcal E_8\), the family
\(\mathsf{RosetteSpec}_8(X_8)\) has no constructor.  The canonical-forms
theorem therefore shows that it has no closed inhabitant.  The extension adds
exactly the payload, event, and aligned origin required by
\textsc{Sig-Certificate}; it adds no reduction rule.

The following table is a finite admission certificate for the staged module.
\begin{table}[ht]
\centering
\small
\begin{tabularx}{\textwidth}{@{}l>{\raggedright\arraybackslash}p{.24\textwidth}X@{}}
\toprule
Entry & Admission rule & Discharged condition\\
\midrule
\(c_8,h_8\) & \textsc{Sig-Primitive} & closed, grade-free,
family-headed form types\\
\(\mathsf{RosetteSpec}_8\) & \textsc{Sig-Family} & form dependency is
ordinary into witnesses\\
\(\mathsf{inlay}_8\) & \textsc{Sig-Gate} & form parameter demand \(0\),
witness credential demand \(1\), protected grade-free result\\
\(d_8\) & \textsc{Sig-Certificate} & closed non-rendering family-headed
payload; fresh common owner\\
\bottomrule
\end{tabularx}
\caption{Admission certificate for the octagonal-rosette gate and its
evidence extension.}
\label{tab:rosette-certificate}
\end{table}

The gate can be postponed behind a higher-order interface, but only at the
positive grade demanded by the binder policy:
\[
 W_8=
 \lambda^1(s:_{\Wit}\Stamped_\Wit(\mathsf{RosetteSpec}_8(X_8))).
 \mathsf{gate}_{\mathsf{inlay}_8}(X_8;s).
\]
The body component of \(s\) is exactly \(1\), so maximal elaboration derives
\[
 \vdash_{\Frm}W_8:
 \AFDepFun^1(s:_{\Wit}\Stamped_\Wit(\mathsf{RosetteSpec}_8(X_8))).
 \Form.
\]
This function type is intentionally not grade-free: \(W_8\) is a suspended
authorization and has not emitted a record.  Applying it to the extension
certificate produces the old protected result type \(\Form\), where the
grade-free theorem applies:
\begin{align*}
 T_8
 &:=W_8\,\mathsf{stamp}(a_8,e_8,o_8)\\
 &\longrightarrow_\beta
 \mathsf{gate}_{\mathsf{inlay}_8}
   (X_8;\mathsf{stamp}(a_8,e_8,o_8))\\
 &\longrightarrow_{\mathsf{record}}^{\ell_8}
 N_8:=\mathsf{node}_{\mathsf{inlay}_8}(X_8;a_8,e_8).
\end{align*}
All fields are constants or the normal pair \(X_8\), so \(N_8=\nf(T_8)\).
Writing \(\xi_8=\unit_{c_8}\oplus\unit_{h_8}\), the first two terms have
effect
\[
 \unit_{\mathsf{inlay}_8}\oplus\xi_8\oplus
 1\otimes(\unit_{a_8}\oplus\unit_{e_8}\oplus\unit_{o_8}),
\]
and the node has the same effect without the \(o_8\)-component.  The residual
trace and support are
\[
 (\ell_8,e_8,e_8)\in\cross^\downarrow(T_8\longrightarrow^*N_8,N_8),
 \qquad e_8\in\events(N_8).
\]
Thus \(\mathcal E_8\) is not term-conservative at \(\Form\): it enables the
non-old normal form \(N_8\).  It satisfies no-silent-enabling exactly, since
the enabling module's qualified event occurs in that form.

For a separate observation layer, extend by
\[
 \mathcal V_8=
 \{\vdash_{\Ren}\AFCanvas_8\;\jtype,
 \quad
 \mathsf{draw}_8:_{\Ren}
   \AFDepFun^0(X:_{\Frm}\Form).\AFCanvas_8\}.
\]
This suffix is rendering-only over the already extended signature.
\Cref{thm:rendering,thm:composition} show that adding it neither creates
another old protected form nor changes equality between existing forms.

Finally, erasure makes the operational loss of the certificate explicit:
\begin{align*}
 |W_8|&=\lambda s.\overline{\mathsf{inlay}}_8(|X_8|;s),\\
 |T_8|&\longrightarrow_\beta
 \overline{\mathsf{inlay}}_8(|X_8|;|a_8|)
 =|N_8|.
\end{align*}
The source record step is the equality on the right, hence a stuttering
erasure step; \(e_8\) and \(o_8\) are absent.  As a concrete finite structural
interpretation, take a vocabulary with binary relations
\(\mathsf{Rim}\) and \(\mathsf{Spoke}\) and set
\[
 \sem{|N_8|}_{8}=
 \bigl(\fin(9),
   \{(i,1+(i\bmod 8))\mid1\le i\le8\},
   \{(0,i)\mid1\le i\le8\}\bigr).
\]
This is the directed rim and eight spokes of a wheel.  Its canonical finite
table records the realized rosette, while the source normal form additionally
records which certificate event authorized it.

\subsection{Installed presentations}
\label{sec:installed-structure}

The book's \PRES{} retains a selected order, seams, coordinates, cuts and
sectors.  A description offered to a constructor is witness data; once
installed it is presentation, and the installation is a gate.  For a selected
order,
\[
 \vdash_{\Wit}\mathsf{OrderSpec}(F)\;\jtype,
 \qquad
 \vdash_{\Prs}\mathsf{Ordered}(F)\;\jtype,
\]
and the presentation gate is
\[
 \mathsf{installOrder}:
 (F:_{\Frm}\Form);\,
 \Stamped_\Wit(\mathsf{OrderSpec}(F))
 \Longrightarrow_{\Prs}\mathsf{Ordered}(F).
\]
The form \(F\) is an ordinary parameter.  It is not recovered by projecting a
presentation-role pair back into the form role.  An artifact consists
externally of the already constructed \(F\) and a term of
\(\mathsf{Ordered}(F)\); the normal form of the presentation records the
installation event without pretending that the event constructed \(F\).

The book's base cut of a wheel is this pattern: offered to the constructor
that points the wheel it is a witness, and carried by the pointed wheel it is
presentation \cite[\S1.9]{Book}.  So is a supplied policy: witness data while
it is supplied, presentation once a realizer carries it (\cref{sec:realizer}).
Coordinates, seams, cuts and sectors use the same indexed pattern.  A
descriptor lives at \(\Wit\), the installed family at \(\Prs\), and
installation is a presentation gate.

\subsection{Interpretation witnesses and the raw/interpreted split}
\label{sec:interpretation-witnesses}

The book reads one raw artform in several ways without changing
it: a palette laid on a clock by recurrence or by division, a template
recognized in a canvas, a connection read through an invariant.  The book's
interpretation witness is a realizer paired with the certificate that the
realizer produces an endpoint from the artform and obeys the law of a
profile, and interpretation is a species of the witness role rather than a
sixth role \cite[Def.~1.9]{Book}.
In \(\RATT\) that architecture is a declaration schema at role \(\Wit\), and
the rules fix exactly what the schema may package.

An \emph{interpretation profile} \(\mathcal C\) over a raw type
\(\mathsf{Art}\), which is \(\Form\) or a declared form-role family,
consists of an endpoint family and two witness families,
\[
\begin{gathered}
 \Theta_{\mathcal C}:(F:_{\Frm}\mathsf{Art})\Rightarrow_{\rho_{\mathcal C}}\Type,
 \qquad
 \mathcal R_{\mathcal C}:
  (F:_{\Frm}\mathsf{Art},\ \theta:_{\rho_{\mathcal C}}\Theta_{\mathcal C}(F))
  \Rightarrow_{\Wit}\Type,\\
 \mathcal L_{\mathcal C}:
  (F:_{\Frm}\mathsf{Art},\ \theta:_{\rho_{\mathcal C}}\Theta_{\mathcal C}(F),\
   \ell:_{\Wit}\mathcal R_{\mathcal C}(F,\theta))
  \Rightarrow_{\Wit}\Type,
\end{gathered}
\]
with \(\rho_{\mathcal C}\ne\Ren\), together with ordinary constructors for
the two witness families.  Every telescope edge into \(\Wit\) has demand
\(0\) in \cref{tab:policy}, so \textsc{Sig-Family} admits both families.
The interpretation witness type is the derived homogeneous sum
\[
 \AFIntWit_{\mathcal C}(F,\theta)
 \coloneqq
 \AFDepPair(\ell:_{\Wit}\mathcal R_{\mathcal C}(F,\theta)).\,
 \mathcal L_{\mathcal C}(F,\theta;\ell),
\]
formed by \textsc{Ty-Sigma} at role \(\Wit\).  Its first projection is the
supplied realizer and its second the law certificate; the source \(F\), the
profile, and the endpoint \(\theta\) remain indices, and the source enters
only through ordinary form edges.  Because both components are
witness-family applications, \((\mathrm{WitnessData})\) gives
\(\mathsf{wdata}(\AFIntWit_{\mathcal C}(F,\theta))\): a closed public
inhabitant \(I=(\ell,c)\) is exactly attestable by \textsc{Sig-Attest}
without a fresh payload constant.

Three features of the core determine what the schema does not contain.
\begin{enumerate}[label=(\roman*),leftmargin=2.5em]
\item There is no internal type of profiles or endpoints, and dependent sums
are homogeneous in role.  The book's semantic fibre, the sum
over the available profiles and their endpoints of
\(\AFIntWit_{\mathcal C}(F,\theta)\), is therefore not a \(\RATT\)
type.  It is the external family of witness types indexed by signature
entries and endpoint terms.  For the same reason the layer pairing a raw
\(F\) with one selected witness is not a type: as for
\(\mathsf{Ordered}(F)\) in \cref{sec:installed-structure}, an artifact with
a selected reading consists externally of the already constructed \(F\) and
a term of \(\AFIntWit_{\mathcal C}(F,\theta)\).  The book's fibre stands
beside \(F\) and selects nothing, and choosing one witness is a further step;
here both are consequences of the universe-free boundary of
\cref{sec:boundary}.
\item An endpoint cannot be a rendering.  A rendering-role index has demand
\(\absent\) into \(\Wit\), so a profile whose endpoints are graphics, scores,
or observed values fails policy validation.  A value read off an artform
enters only as a closed literal or meta-level finite index, which carries no
declaration atom.
\item Forming a witness installs nothing.  A selected inhabitant becomes
an installed presentation only through a presentation gate
\[
\begin{gathered}
 \mathsf{install}_{\mathcal C}:
 (F:_{\Frm}\mathsf{Art},\ \theta:_{\rho_{\mathcal C}}\Theta_{\mathcal C}(F));\\
 \Stamped_\Wit(\AFIntWit_{\mathcal C}(F,\theta))
 \Longrightarrow_{\Prs}\mathsf{Interp}_{\mathcal C}(F,\theta)
 \quad[\mathsf{node}_{\mathsf{install}_{\mathcal C}}],
\end{gathered}
\]
where \(\mathsf{Interp}_{\mathcal C}\) is a presentation-role family over the
same telescope.  \textsc{Sig-Gate} requires
\(\rho_{\mathcal C}\in\{\Frm,\Prs\}\) by (G1), since a witness- or
script-role parameter has demand \(1\) into \(\Prs\); the payload type is
grade-free; and the result does not depend on the payload.  A certificate
bundle cannot deliver the credential, because \textsc{Sig-Certificate}
requires a family-headed payload type and
\(\AFIntWit_{\mathcal C}(F,\theta)\) is a sum.  The route into the
gate is exact attestation of an already checked pair.
\end{enumerate}

\begin{proposition}[Interpretation profiles are conservative]
\label{prop:profile-conservative}
Let \(\mathcal M_{\mathcal C}\) add, over a signature in which
\(\mathsf{Art}\) and \(\Theta_{\mathcal C}\) are old, exactly the two
witness families \(\mathcal R_{\mathcal C}\), \(\mathcal L_{\mathcal C}\)
and ordinary constructors headed by them.  Then \(\mathcal M_{\mathcal C}\)
is a fresh-witness-family module, and by \cref{thm:fresh} it is
term-conservative and equality-reflecting at every old closed grade-free
type of role \(\Frm\) or \(\Prs\).  In particular a new way of reading
\(F\) creates no raw artform, identifies none, and installs no presentation.
\end{proposition}

\begin{proof}
The suffix consists of witness-role families and constructors headed by
them, which is the fresh-witness-family class of \cref{sec:modules}; the
rest is \cref{thm:fresh}.
\end{proof}

The hypothesis that \(\Theta_{\mathcal C}\) is old is exact.  A profile that
also declares a new form- or presentation-role endpoint family leaves the
class, and no conservativity claim is exported for it.

\begin{proposition}[Attested installation is traceable]
\label{prop:attested-installation}
Let \(I\) be a closed public term of
\(\AFIntWit_{\mathcal C}(F,\theta)\) in a signature containing
\(\mathcal M_{\mathcal C}\) and \(\mathsf{install}_{\mathcal C}\), and let
\(\mathcal E_q\) be the exact attestation \(q@I=(e_q,o_q)\).
Then \(\mathcal E_q\) is an evidence-inhabitant module, and the public term
\[
 \mathsf{gate}_{\mathsf{install}_{\mathcal C}}
 (F,\theta;\mathsf{stamp}(I,e_q,o_q))
\]
contracts to
\(\mathsf{node}_{\mathsf{install}_{\mathcal C}}(F,\theta;I,e_q)\), with the
exact pair retained and \(e_q\) in its event support; and, by
\cref{thm:no-silent}, every closed public term of an old grade-free
protected type whose normal form is not old contains an event owned by
\(\mathcal E_q\).
\end{proposition}

\begin{proof}
The attestation is admitted by \textsc{Sig-Attest}, whose premises were
verified above, and it adds neither family, primitive, nor gate, so the
module is evidence-inhabitant.  The contraction is one instance of the
record rule, and \cref{lem:record} carries \(e_q\) to the normal form.
The last claim is \cref{thm:no-silent}.
\end{proof}

We now give one finite instance, worked out against the printed rules in
the same way as \cref{sec:rosette-module}.  Extend the Artforms signature with
an endpoint family and closed constants
\[
\begin{gathered}
 \mathsf{Tap}:(F:_{\Frm}\Form)\Rightarrow_{\Frm}\Type,
 \qquad
 c_Q:_{\Frm}\AFFinCode,\quad
 h_Q:_{\Frm}\Hist(c_Q),\quad
 Q\coloneqq(c_Q,h_Q),\\
 \theta_{\mathrm{rec}},\theta_{\mathrm{div}}:_{\Frm}\mathsf{Tap}(Q),
\end{gathered}
\]
where \(Q\) is intended as a palette, \(\mathsf{Tap}(Q)\) as the tapestries
laid from it, and the two constants as the tapestries obtained by laying the
palette by recurrence and by division.  A laid tapestry is a form, as the
book's tapestries are, so the endpoints are form-role data.  The recurrence profile module
\(\mathcal M_{\mathrm{rec}}\) consists of
\[
\begin{gathered}
 \mathsf{Lay}_{\mathrm{rec}}:
  (F:_{\Frm}\Form,\ \theta:_{\Frm}\mathsf{Tap}(F))\Rightarrow_{\Wit}\Type,\\
 \mathsf{Obeys}_{\mathrm{rec}}:
  (F:_{\Frm}\Form,\ \theta:_{\Frm}\mathsf{Tap}(F),\
   \ell:_{\Wit}\mathsf{Lay}_{\mathrm{rec}}(F,\theta))\Rightarrow_{\Wit}\Type,\\
 \ell_{\mathrm{rec}}:_{\Wit}\mathsf{Lay}_{\mathrm{rec}}(Q,\theta_{\mathrm{rec}}),
 \qquad
 c_{\mathrm{rec}}:_{\Wit}
  \mathsf{Obeys}_{\mathrm{rec}}(Q,\theta_{\mathrm{rec}},\ell_{\mathrm{rec}}),
\end{gathered}
\]
whose intended fields are the rule that advances every row by one phase per
tick and the check that its evaluation is \(\theta_{\mathrm{rec}}\);
\(\mathcal M_{\mathrm{div}}\) is the same with \(\mathrm{div}\) and the rule
that holds a phase for a whole bar.  Both modules are fresh-witness-family,
so \cref{prop:profile-conservative} applies: after either or both, the
closed form terms \(c_Q\), \(h_Q\), \(Q\) and every old closed term of
type \(\Form\) are the same old terms with the same equalities.  The two
witnesses
\[
 I_{\mathrm{rec}}\coloneqq(\ell_{\mathrm{rec}},c_{\mathrm{rec}})
 :\AFIntWit_{\mathrm{rec}}(Q,\theta_{\mathrm{rec}}),
 \qquad
 I_{\mathrm{div}}\coloneqq(\ell_{\mathrm{div}},c_{\mathrm{div}})
 :\AFIntWit_{\mathrm{div}}(Q,\theta_{\mathrm{div}})
\]
are closed normal terms of two witness-role types over the one form
constant \(Q\).  Nothing relates \(\theta_{\mathrm{rec}}\) to
\(\theta_{\mathrm{div}}\): the core has no identity type and no user
equations, and by \cref{thm:canonical-forms} two distinct constants are
distinct normal forms.  ``The same palette, laid twice'' is thus literally
two readings beside one unchanged \(Q\).

Installation of the recurrence reading is the gate module
\[
\begin{gathered}
 \mathsf{Laid}_{\mathrm{rec}}:
  (F:_{\Frm}\Form,\ \theta:_{\Frm}\mathsf{Tap}(F))\Rightarrow_{\Prs}\Type,\\
 \mathsf{installLaying}_{\mathrm{rec}}:
 (F:_{\Frm}\Form,\ \theta:_{\Frm}\mathsf{Tap}(F));\\
 \Stamped_\Wit(\AFIntWit_{\mathrm{rec}}(F,\theta))
 \Longrightarrow_{\Prs}\mathsf{Laid}_{\mathrm{rec}}(F,\theta)
 \quad[\mathsf{node}_{\mathsf{installLaying}_{\mathrm{rec}}}],
\end{gathered}
\]
and the evidence-inhabitant extension is the single exact attestation
\(\mathcal E_{\mathrm{rec}}=\{q@I_{\mathrm{rec}}=(e_{\mathrm{rec}},o_{\mathrm{rec}})\}\).
\Cref{tab:interpretation-certificate} is the admission certificate.

\begin{table}[ht]
\centering
\small
\begin{tabularx}{\textwidth}{@{}l>{\raggedright\arraybackslash}p{.24\textwidth}X@{}}
\toprule
Entry & Admission rule & Discharged condition\\
\midrule
\(\mathsf{Tap}\); \(\theta_{\mathrm{rec}},\theta_{\mathrm{div}}\) &
\textsc{Sig-Family}; \textsc{Sig-Primitive} & form dependency ordinary
into form; closed, grade-free, family-headed form types\\
\(\mathsf{Lay}_{\mathrm{rec}}\), \(\mathsf{Obeys}_{\mathrm{rec}}\) &
\textsc{Sig-Family} & form and witness dependencies ordinary into
witnesses\\
\(\ell_{\mathrm{rec}}\), \(c_{\mathrm{rec}}\) & \textsc{Sig-Primitive} &
closed, grade-free, headed by the new witness families\\
\(\AFIntWit_{\mathrm{rec}}(Q,\theta_{\mathrm{rec}})\) &
\textsc{Ty-Sigma} & homogeneous at \(\Wit\); witness data\\
\(\mathsf{Laid}_{\mathrm{rec}}\); \(\mathsf{installLaying}_{\mathrm{rec}}\) &
\textsc{Sig-Family}; \textsc{Sig-Gate} & form parameters demand \(0\)
into presentation; witness credential demand \(1\); protected grade-free result
independent of the payload\\
\(q@I_{\mathrm{rec}}\) & \textsc{Sig-Attest} & closed public witness-data
payload; grade-free; no payload constant; fresh common owner\\
\bottomrule
\end{tabularx}
\caption{Admission certificate for the recurrence interpretation profile,
its installation gate, and its exact attestation.}
\label{tab:interpretation-certificate}
\end{table}

The public term and its record step are
\begin{align*}
 T_{\mathrm{rec}}&\coloneqq
 \mathsf{gate}_{\mathsf{installLaying}_{\mathrm{rec}}}
  (Q,\theta_{\mathrm{rec}};
   \mathsf{stamp}(I_{\mathrm{rec}},e_{\mathrm{rec}},o_{\mathrm{rec}}))\\
 &\longrightarrow_{\mathsf{record}}^{\ell_{\mathrm{rec}}}
 N_{\mathrm{rec}}\coloneqq
 \mathsf{node}_{\mathsf{installLaying}_{\mathrm{rec}}}
  (Q,\theta_{\mathrm{rec}};I_{\mathrm{rec}},e_{\mathrm{rec}}).
\end{align*}
Writing \(\xi_Q=\unit_{c_Q}\oplus\unit_{h_Q}\), the effect of
\(T_{\mathrm{rec}}\) is
\[
 \unit_{\mathsf{installLaying}_{\mathrm{rec}}}\oplus\xi_Q\oplus
 \unit_{\theta_{\mathrm{rec}}}\oplus
 1\otimes(\unit_{\ell_{\mathrm{rec}}}\oplus\unit_{c_{\mathrm{rec}}}
 \oplus\unit_{e_{\mathrm{rec}}}\oplus\unit_{o_{\mathrm{rec}}}),
\]
and the node has the same effect without the \(o_{\mathrm{rec}}\)-component:
subject reduction with effect improvement.  All fields are constants or
normal pairs, so \(N_{\mathrm{rec}}=\nf(T_{\mathrm{rec}})\),
\(e_{\mathrm{rec}}\in\events(N_{\mathrm{rec}})\), and
\((\ell_{\mathrm{rec}},e_{\mathrm{rec}},e_{\mathrm{rec}})
 \in\cross^\downarrow(T_{\mathrm{rec}}\longrightarrow N_{\mathrm{rec}},
 N_{\mathrm{rec}})\).
The exact pair \(I_{\mathrm{rec}}\) is retained in the node, and no
projection equation was needed to recover its realizer.  A second
attestation \(q'@I_{\mathrm{rec}}=(e'_{\mathrm{rec}},o'_{\mathrm{rec}})\)
installs the same reading with a different event: the two nodes are
distinct source terms with the same erasure
\[
 |N_{\mathrm{rec}}|=
 \overline{\mathsf{installLaying}}_{\mathrm{rec}}
 (|Q|,|\theta_{\mathrm{rec}}|;|I_{\mathrm{rec}}|),
\]
which forgets \(e_{\mathrm{rec}}\) and \(o_{\mathrm{rec}}\) as
\cref{thm:erasure} prescribes.  Installing the division reading is the same
derivation with \(\mathrm{div}\), and its node carries \(e_{\mathrm{div}}\).
A finite structural interpretation may then denote \(|N_{\mathrm{rec}}|\)
and \(|N_{\mathrm{div}}|\) by two different tables over the same carrier,
while only the source normal forms record which reading and which
attestation authorized each.

\subsection{Frames, states and scripts}\label{sec:frames-calculus}

A frame of the book is a carrier, its carried structure, the space of legal
states, the allowed observations and the allowed moves.  Its first two parts
are the forms and installed presentations above.  For a presentation \(A\)
declare
\[
\begin{gathered}
 \vdash_{\Wit}\mathsf{Table}(A)\;\jtype,
 \qquad
 \mathsf{Legal}_A:(s:_{\Wit}\mathsf{Table}(A))\Rightarrow_{\Wit}\Type,\\
 \State(A)\coloneqq
 \AFDepPair(s:_{\Wit}\mathsf{Table}(A)).\,\mathsf{Legal}_A(s).
\end{gathered}
\]
A table is a choice of a value at every position of \(A\), and a state is a
table with the evidence that it is legal.  The sum is homogeneous at
\(\Wit\), so \textsc{Ty-Sigma} forms it: the book's formation rule for the
state space is a witness-role sum, as \cref{sec:states} reads it.  The
presentation \(A\) enters only as an index, through an ordinary edge from
\(\Prs\) into \(\Wit\).

For presentations \(A\) and \(B\) declare
\[
\begin{gathered}
 \vdash_{\Scr}\Script(A,B)\;\jtype,
 \qquad
 \mathsf{id}_A:_{\Scr}\Script(A,A),\\
 \mathsf{comp}:_{\Scr}
 \AFDepFun^0(f:_{\Scr}\Script(A,B)).\,
 \AFDepFun^0(g:_{\Scr}\Script(B,C)).\,\Script(A,C),\\
 \mathsf{act}_{\mathsf T}:_{\Wit}
 \AFDepFun^0(f:_{\Scr}\Script(A,B)).\,
 \AFDepFun^0(t:_{\Wit}\mathsf{Table}(A)).\,\mathsf{Table}(B),\\
 \mathsf{act}_{\mathsf L}:_{\Wit}
 \AFDepFun^0(f:_{\Scr}\Script(A,B)).\,
 \AFDepFun^0(t:_{\Wit}\mathsf{Table}(A)).\,
 \AFDepFun^0(\ell:_{\Wit}\mathsf{Legal}_A(t)).\,
 \mathsf{Legal}_B(\mathsf{act}_{\mathsf T}\,f\,t).
\end{gathered}
\]
Script generators are further primitives at \(\Scr\).  Composition is typed
at the middle presentation, as in the book.  The action of a script on a
state is the derived term
\[
 \mathsf{act}\,f\,p\coloneqq
 \bigl(\mathsf{act}_{\mathsf T}\,f\,(\fst p),\
 \mathsf{act}_{\mathsf L}\,f\,(\fst p)\,(\snd p)\bigr):\State(B),
\]
which sends a state of the source to a state of the target and keeps it
legal: a script is total on legal states.  Its result is witness data, since
the codomain decides.  Every edge in these declarations has demand \(0\),
since \(D(\Scr,\Wit)=D(\Wit,\Wit)=0\), so they are ordinary primitives and
need no gate.  A move changes the instance and not the frame.

The core treats script equality as judgmental equality of the freely generated
syntax only.  Two scripts may have the same action and remain two scripts, as
the book requires: that their actions agree is itself witness data, an
inhabitant of a family
\(
 \mathsf{SameAction}(f,g)
\)
at \(\Wit\) indexed by the two scripts.  Fusion, commutation or cancellation
equations belong to a separate typed setoid, after a proof that their
congruence reflects old equality; they do not participate in dependent
conversion.

Evidence that a move is allowed is a witness family indexed by the script and
the state,
\[
 \mathsf{Applicable}:(f:_{\Scr}\Script(A,B),\ s:_{\Wit}\State(A))
 \Rightarrow_{\Wit}\Type,
\]
admitted because both edges are ordinary into \(\Wit\).  Requirement (R4)
appears here twice: the certificate names the move, and a guarded action may
take its certificate.  The step from witnesses and scripts back into the
frame is gated.  A position that becomes the fixed start of a puzzle, or a
history that becomes a game score, enters a form or a presentation only
through a gate, and its normal form records the event that authorized it.  For
a score,
\[
 \mathsf{note}:
 (T:_{\Frm}\mathsf{Score});\,
 \Stamped_\Wit(\mathsf{Move})
 \Longrightarrow_{\Frm}\mathsf{Score},
\]
where \(\mathsf{Move}\) is a closed witness family of move records; each
contraction retains the move and its event, so a score in normal form is the
list of the moves made with their events.  This is the book's run, a
trajectory together with the trace of the moves made, with the trace as the
record of the gates.

\subsection{One capture in the calculus}\label{sec:capture}

The capture of \cref{tab:chess} is a closed derivation in this signature.  Let
\(\mathsf{Board}\coloneqq(c_B,h_B):\Form\) be the board, built from constants
as \(X_8\) is in \cref{sec:rosette-module};
\(A:_{\Prs}\mathsf{Addr}(\mathsf{Board})\) the addressing;
\(s_0\coloneqq(t_0,\ell_0):\State(A)\) the position, with
\(t_0:_{\Wit}\mathsf{Table}(A)\) and \(\ell_0:_{\Wit}\mathsf{Legal}_A(t_0)\);
and \(\mathsf{nxe5}:_{\Scr}\Script(A,A)\) the capture.  The squares White bears
on are a derived state,
\[
 \mathsf{attack}:_{\Wit}
 \AFDepFun^0(s:_{\Wit}\State(A)).\,\mathsf{Table}_{\mathsf{Bool}}(A),
\]
a value on every square, and the certificate of the capture is a term of
\(\mathsf{Applicable}(\mathsf{nxe5},s_0)\) whose derivation may use
\(\mathsf{attack}\,s_0\).  The position after the capture is
\(\mathsf{act}\,\mathsf{nxe5}\,s_0\); its effect has the component \(0\) at
\(\mathsf{act}_{\mathsf T}\), \(\mathsf{act}_{\mathsf L}\), \(\mathsf{nxe5}\),
\(t_0\) and \(\ell_0\), and at the declarations of \(A\) and the board.  The engine's evaluation is
\[
 \mathsf{eval}:_{\Ren}\AFDepFun^0(s:_{\Wit}\State(A)).\,\mathsf{Num},
 \qquad
 \vdash_{\Ren}\mathsf{Num}\;\jtype .
\]
A certificate family indexed by \(\mathsf{eval}\,s_0\) cannot be formed: its
formation edge from a rendering into \(\Wit\) has demand \(\absent\), so
\textsc{Sig-Family} rejects it.  A rule that consults a reading is not merely
discouraged; it is not a well-formed declaration.  The derivations' effects
are the grades of the construction of \cref{tab:chess}, item for item, and
\cref{prop:chess} is the statement that this signature is admitted while its
two variants are not.

\subsection{Rendering firewall}

A rendering is a value at a rendering-role type, and a readout is a script or
function whose values are renderings.  Declare
\[
  \vdash_{\Ren}\Graphic\;\jtype,
  \qquad
  r:_{\Ren}\AFDepFun^0(s:_{\Wit}\State(A)).\Graphic.
\]
The dependency of the drawing on the state is ordinary.  Conversely, a
function type with a graphic domain and a non-rendering result fails policy
validation unless the domain edge is \(\absent\); no primitive may declare
such an edge, stamps cannot contain rendering payloads, and gates cannot
target rendering data.  Thus replacing the renderer cannot alter an
applicability derivation, which is \cref{thm:rendering} for this signature and
\cref{thm:firewall-values} for its values.

A policy that chooses among applicable moves is witness data while it is
supplied.  It may depend on states and certificates, but not on a rendering;
a scheduler that inspects a drawing to choose a move is outside the calculus.
Installed in a realizer, the policy becomes presentation through a gate, as
in \cref{sec:installed-structure}.

The protected observation types used above---\(\Form\), \(\mathsf{Score}\),
\(\mathsf{Ordered}(F)\), and every admitted gate result---are grade-free.
Consequently the adequacy and module theorems apply to completed forms,
scores and installed presentations.  A higher-order scheduler interface whose
type itself contains grade \(1\) or \(\absent\) is intentionally outside that
observation boundary until it is applied to a closed implementation.

\subsection{Finite structural semantics}

For a fixed finite relational vocabulary \(L\), an Artforms signature may
supply a total interpretation
\[
  \sem{-}_L:
  |\Form|\longrightarrow\mathsf{FinStr}_L
\]
on erased normal forms.  Canonical labeling over all bijections to
\(\fin(n)\) decides isomorphism of the resulting finite structures.  The map
factors as
\[
 \text{RATT term}
 \longrightarrow
 \text{provenance normal form}
 \xrightarrow{\erase}
 \text{finite structure}
 \longrightarrow
 \text{canonical table}.
\]
The first arrow preserves gate events; the second intentionally forgets them.
Different construction histories may therefore denote the same finite
structure without becoming judgmentally equal.

\section{Relationship to existing work}\label{sec:related}

\paragraph{The textbook.}
The five roles, their readings, the policy in prose and its key idea that
nothing may depend on a reading are the book's \cite[\S1.9]{Book}.  So are the
frame and its instance, the four useful words that are not roles, the rules of
inference read as a second discipline beside the roles \cite[\S1.11]{Book},
and the count of histories that holds moves and not proofs
\cite[\S7.4]{Book}.  This paper adds the levels, the uniqueness and path
closure of the policy, the separation theorem, the decision that a state is
witness data, and the proofs.  The book does not cite this paper; the
repository that holds the Lean proofs of the book's theorems also holds this
paper and its Lean library \cite{Repo}.

\paragraph{Graded and multimodal dependent type theory.}
Multimodal dependent type theory parameterizes syntax by modes, modalities,
and transformations and proves uniform canonicity and normalization results
\cite{MTT,MTTNorm}.  Graded modal dependent type theory tracks variable use in
terms and types and establishes normalization, decidable equality, and erasure
properties \cite{GradedDTT}.  \(\RATT\) does not claim roles or graded
substitution as new in isolation.  Its grade algebra measures the least
mediation over occurrence paths, and its generative restriction ties the only
positive primitive grade to a record-emitting reduction rule.

\paragraph{Phase separation and noninterference.}
Modal type theories connect phase distinctions with information-flow
noninterference \cite{KavvosIF}.  Type Theory with Erasure gives a structural
phase distinction and conservativity over Martin-L\"of type theory
\cite{TypeTheoryErasure}.  Abstraction Functions as Types derives modularity
from a phase distinction in univalent dependent type theory
\cite{AbstractionFunctions}.  RATT's rendering theorem is a small generative
noninterference result, while its evidence theorem deliberately permits a
declassification path and requires an intrinsic record.

\paragraph{Dependency security and declassification.}
The Dependency Core Calculus internalizes protection labels in types and proves
noninterference properties for dependency \cite{CoreDependency}; systematic
accounts of declassification distinguish what information is released, by
whom, where, and when \cite{Declassification}.  The four levels of
\cref{sec:levels} order the roles as a lattice of labels orders data, and the
firewall is a noninterference theorem.  But \(\RATT\) does not claim a
confidentiality theorem, and its roles are not security levels: a gated
crossing is an installation that is recorded, not a release of a secret.  Its narrower
connection is controlled release into a protected constructor: the release
edge is generative, and a contributing release leaves a proof-relevant record
in the normal form.  The grade-free boundary in
\cref{rem:gradefree-boundary} is necessary because an unapplied higher-order
interface can postpone such a release without yet producing a record.

\paragraph{Provenance.}
Tracing calculi for functional programs use replayable traces and support
multiple provenance views \cite{CoreProvenance}.  RATT differs by making a
record constructor part of the normal form of a protected dependent term and
by proving that the static grade-\(1\) path and dynamic event arise from the
same gate declaration.

\paragraph{Modular signatures.}
Generalized algebraic theories and categories with families provide standard
accounts of dependent signatures, while modular presentations of type theory
study when subtheory derivations remain in a smaller signature
\cite{DybjerCwF,ModularTT}.  RATT selects a deliberately generative fragment
whose extension classes have different, syntactically checkable theorems.

\paragraph{Mechanization.}
Graded modal dependent type theory with a universe and erasure has been
formalized in Agda \cite{GradedDTT}.  Here the first part is formalized in
Lean~4 with Mathlib \cite{Lean4,Mathlib}; the calculus of the second part is
not yet formalized, and \cref{sec:status} says what that would take.

\section{Theorem status and limitations}\label{sec:status}

\begin{table}[ht]
\centering
\small
\begin{tabularx}{\textwidth}{@{}>{\raggedright\arraybackslash}p{.30\textwidth}
 >{\raggedright\arraybackslash}p{.17\textwidth}X@{}}
\toprule
Result & Status & Exact boundary\\
\midrule
Levels, uniqueness and path closure of the policy & Lean &
\cref{thm:levels,thm:forced,thm:path-closed}\\
Policy separation & Lean & \cref{thm:policy-separation}\\
Local to global; firewall; mediation; records & Lean & constructions with
graded dependencies (\cref{sec:constructions})\\
Firewall, frame stability, presentation independence for values & Lean &
acyclic constructions; evaluations that read live dependencies only\\
Three counts of histories; separation of \(\Wit\) and \(\Scr\) & Lean & finite
laws (\cref{sec:histories})\\
Capture, realizer, walker & Lean & \cref{sec:examples}\\
\midrule
Principal effects & paper proof & finite min--max effects; explicit grades\\
Substitution & paper proof & types and terms; path-closed policy\\
Subject reduction & paper proof & stable lambda bounds; effects may improve\\
Strong normalization & paper proof & skeleton-indexed reducibility; universe-free core\\
Confluence & paper proof & complete local-peak analysis plus Newman\\
Canonical forms & paper proof & explicit normal/neutral grammar; qualified credentials\\
Decidable conversion & proved corollary & normalization plus confluence\\
Path mediation & paper proof & effect paths; no claim that syntax is erased\\
Grade-free mediation & paper proof & public atoms; closed grade-free protected result\\
Static--dynamic gate adequacy & paper proof & closed public terms; residual node ancestry\\
Residual record adequacy & paper proof & final fields of surviving node residuals\\
Extension support & paper proof & public reachability; old closed grade-free result types\\
Rendering and fresh-family conservativity & proved corollaries & exact boundaries above\\
Evidence traceability & proved corollary & grade-free protected types;
old-family bundles or exact attestations of old witness data\\
Erasure simulation & paper proof & event-forgetting role-free target\\
Rosette, interpretation and capture modules & derivation on paper &
\cref{sec:artforms}\\
\midrule
Mechanization of the calculus & open & \cref{sec:core,sec:reduction,sec:metatheory,sec:adequacy,sec:modules,sec:erasure} are not machine checked\\
Role-indexed universes & open extension & no internal universe in the core\\
Quotients and forced values & open extension & require separate computation or staging\\
\bottomrule
\end{tabularx}
\caption{Formal claims and their validation boundaries.  \emph{Lean} means
checked by Lean's kernel, with the names of \cref{tab:lean}.}
\label{tab:status}
\end{table}

Here ``paper proof'' means that the argument is carried out for the displayed
rules, including the induction measure or case analysis needed by every
nonstandard constructor.  A ``proved corollary'' invokes those paper proofs.
Neither denotes machine verification.

The principal remaining task is to mechanize the calculus.  The first part
already supplies its algebra: the grades, the policy and its path closure,
and the chains of dependencies that \cref{lem:path-characterization} reads off
a derivation.  The next step is a simply typed fragment with stamps, gates
and records, where substitution, preservation and the record theorems are
modest; the dependent core with its reducibility proof of strong
normalization is a larger project.  A second task is to add role-indexed
universes without reintroducing fresh-type packaging into protected old
types.  A third is to give certified inductive families and recursion through
a concrete host checker rather than broad proof-carrying rewrite
certificates.

The calculus is intentionally intensional.  It does not identify different
histories with the same erased structure, and it does not add an extensional
quotient of scripts.  Those choices keep definitional equality decidable and
make the provenance theorem literal rather than modulo an unspecified
congruence; they are also the book's role discipline, under which equal
derived values do not identify constructions, presentations, witnesses or
transformation histories.

\section{Conclusion}\label{sec:conclusion}

The book sorts every datum into five roles and lets one rule carry most of
the weight.  The first part of this paper shows that the policy behind that
rule is not a list of conventions.  It is generated by four levels---the
carrier, its carried structure, the instance, and the display---and four
requirements the book states force it.  It is path closed, so a construction
can be checked one dependency at a time, and on every valid construction the
firewall, recorded mediation, frame stability and presentation independence
hold, for chains and for values.

The policy separates nine of the ten pairs of roles.  The tenth pair,
witnesses and scripts, shares the instance level because the book needs
certificates about moves and moves that take certificates.  Their rules
separate them, and the separation is measurable: a history is a list of
applications, a certificate only admits one, and counting either way
changes the book's counts.  A state is a legal choice of values and so
witness data, and every datum has one of the five roles.

The second part puts the roles into the judgement of a dependent type
theory.  Its min--max effect computes the least mediation on every path, a
registered gate is the only source of a gated dependency, and its contraction
leaves a record node with a module-owned event, so the static path and the
dynamic provenance come from one rule.  Renderings and fresh witness
vocabularies are conservative, new evidence enables construction only on the
record, and the book's frames, states, scripts, interpretation witnesses and
renderings are its declarations.  The first part is checked in Lean; the
second is the next to be.

\appendix

\section{Complete declarative rule set}\label{app:rules}

This appendix makes explicit the administrative rules used in the main text.
There are no additional typing or reduction rules.

\subsection{Raw and accepted judgments}

First erase every policy-validity side condition and derive a \emph{raw}
judgment, written \(\vdash^{0}_{\rho}\).  Its effect is determined by the equations in
\cref{sec:type-rules,sec:term-rules}.  An accepted judgment is defined by
\[
\frac{
 \AFDeclSignature;\Gamma\vdash^{0}_{\rho} J\triangleright\eps
 \qquad \valid_\rho(\eps)
}{
 \AFDeclSignature;\Gamma\vdash_{\rho} J\triangleright\eps
}.
\tag{Accept}
\]
Every premise judgment in a rule is accepted, so illegal dependencies cannot
be hidden in an intermediate type or term.  In addition, every term judgment
\(\vdash_{\rho}t:A\) has an accepted formation derivation
\(\vdash_{\rho}A\;\jtype\); this is an implicit premise of each term rule below and
ensures that grading the term never bypasses a dependency in its type.

The empty and extension context rules are
\[
\frac{\AFDeclSignature\;\wf}{\AFDeclSignature\vdash\cdot\;\ctx}
\tag{Ctx-Empty}
\]
\[
\frac{
 \AFDeclSignature\vdash\Gamma\;\ctx
 \quad
 \AFDeclSignature;\Gamma\vdash_{\rho} A\;\jtype\triangleright\alpha
 \quad x\notin\Gamma
}{
 \AFDeclSignature\vdash\Gamma,x:_{\rho}A\;\ctx
}.
\tag{Ctx-Extend}
\]

A telescope judgment is defined sequentially:
\[
\frac{\AFDeclSignature\vdash\Gamma\;\ctx}
{\AFDeclSignature;\Gamma\vdash\cdot\;\mathsf{tel}\triangleright\absent}
\tag{Tel-Empty}
\]
and
\[
\frac{
 \AFDeclSignature;\Gamma\vdash\Delta\;\mathsf{tel}\triangleright\delta
 \quad
 \AFDeclSignature;\Gamma,\Delta\vdash_{\rho} A\;\jtype\triangleright\alpha
 \quad x\notin\Gamma,\Delta
}{
 \AFDeclSignature;\Gamma\vdash\Delta,x:_{\rho}A\;\mathsf{tel}
 \triangleright\delta\oplus(\alpha\setminus\mathsf{dom}(\Delta))
}.
\tag{Tel-Extend}
\]
Simultaneous term spines are checked from left to right:
\[
\frac{
 \AFDeclSignature;\Gamma\vdash_{\rho_i} t_i:
 A_i[\vec t_{<i}/\vec x_{<i}]\triangleright\eps_i
 \quad(1\le i\le n)
}{
 \AFDeclSignature;\Gamma\vdash\vec t:
 (x_1:_{\rho_1}A_1,\ldots,x_n:_{\rho_n}A_n)
 \triangleright\bigoplus_i\eps_i
}.
\tag{Spine}
\]

\subsection{Signature formation}

The initial signature contains only the built-in formers
\(\Event\), \(\Origin\), and \(\Stamped\):
\[
  \frac{}{\varnothing\;\wf}.
\tag{Sig-Empty}
\]
For each extension rule below, all introduced names are fresh.

\[
\frac{
 \AFDeclSignature\;\wf
 \quad
 \AFDeclSignature;\cdot\vdash\Delta\;\mathsf{tel}
 \quad a\notin\AFDeclSignature
}{
 \AFDeclSignature,a:(\Delta)\Rightarrow_{\rho}\Type\;\wf
}.
\tag{Sig-Family}
\]

Use \(\mathsf{gradefree}(A)\) as defined in \cref{sec:core}, and let
\(\mathsf{ordinaryHead}(A)\) mean that, after removing outer
\(\AFDepFun\)-binders, the result head is a declared family symbol.  Primitive
extension is
\[
\frac{
 \AFDeclSignature\;\wf
 \quad \AFDeclSignature;\cdot\vdash_{\rho} A\;\jtype
 \quad \mathsf{gradefree}(A)
 \quad \mathsf{ordinaryHead}(A)
}{
 \AFDeclSignature,c:_{\rho}A\;\wf
}.
\tag{Sig-Primitive}
\]
The head restriction reserves credentials for the next rule.

\[
\frac{
 \AFDeclSignature\;\wf
 \quad \AFDeclSignature;\cdot\vdash_{\sigma} A\;\jtype
 \quad \sigma\ne\Ren
 \quad \mathsf{gradefree}(A)
 \quad \mathsf{ordinaryHead}(A)
}{
 \AFDeclSignature,
 a_d:_{\sigma}A,
 e_d:_{\Wit}\Event,
 o_d:_{\Wit}\Origin_A(e_d,a_d)
 \;\wf
}.
\tag{Sig-Certificate}
\]
The conclusion is one atomic bundle: a later extension cannot introduce an
origin constant for an event it does not own.

\[
\frac{
 \AFDeclSignature\;\wf
 \quad \AFDeclSignature;\cdot\vdash_{\Wit} t:A
 \quad \mathsf{public}(t)
 \quad \mathsf{gradefree}(A)
 \quad \mathsf{wdata}(A)
}{
 \AFDeclSignature,
 e_q:_{\Wit}\Event,
 o_q:_{\Wit}\Origin_A(e_q,t)
 \;\wf
}.
\tag{Sig-Attest}
\]
The empty context makes \(t\) closed.  This rule introduces no payload
constant and no reduction rule; it records authority for the exact public
witness-data term already checked in \(\AFDeclSignature\).

\begin{remark}[Why certificate heads remain restricted]
If \textsc{Sig-Certificate} admitted a payload type
\(\Stamped_{\Wit}(W)\), its fresh payload constant would be a closed neutral
normal term of stamped type.  An old gate applied to it would be stuck,
contradicting canonical credentials and record progress.  Exact attestation
does not create such a neutral payload, so it does not justify weakening the
certificate rule.
\end{remark}

For gate formation, suppose
\[
 \AFDeclSignature;\cdot\vdash\Delta\;\mathsf{tel},
 \quad
 \AFDeclSignature;\Delta\vdash_{\sigma} A\;\jtype,
 \quad
 \AFDeclSignature;\Delta\vdash_{\rho} B\;\jtype,
\]
where \(\rho\in\{\Frm,\Prs\}\), \(\sigma\ne\Ren\), each declaration role in
\(\Delta\) has demand \(0\) into \(\rho\), and the free variables of \(B\)
are contained in \(\Delta\), and \(\mathsf{gradefree}(B)\).  Then
\[
\frac{\text{the preceding premises}\qquad
      \mathsf{gradefree}(B)\qquad g\text{ fresh}}
{\AFDeclSignature,
 g:\Delta;\Stamped_\sigma(A)\Rightarrow_{\rho} B
 [\mathsf{node}_g]\;\wf}.
\tag{Sig-Gate}
\]
This rule introduces the public atom \(g\) at role \(\rho\), its public gate term former,
and its private node former simultaneously.  It introduces no general
positive-grade constant.

\subsection{Formation and typing summary}

The complete type-former list is:
\[
\begin{array}{c|c|c}
\text{former} & \text{result role} & \text{effect}\\
\hline
a(\vec t) & \rho_a & \unit_a\oplus\bigoplus_i\eps_i\\
\Event & \Wit & \absent\\
\Origin_A(e,t) & \Wit & \alpha\oplus\eta\oplus\theta\\
\Stamped_\sigma(A) & \Wit & \alpha\\
\AFDepFun^q(x:_{\sigma}A).B & \rho &
 q\otimes\alpha\oplus(\beta\setminus x)\\
\AFDepPair(x:_{\rho}A).B & \rho &
 \alpha\oplus(\beta\setminus x)
\end{array}
\]
with the premises and restrictions of
\textsc{Ty-Fam}, \textsc{Ty-Event}, \textsc{Ty-Origin},
\textsc{Ty-Stamp}, \textsc{Ty-Pi}, and \textsc{Ty-Sigma}.  There is no
universe, heterogeneous sum, quotient, identity type, or other implicit type
former in \(\RATT\).  The \textsc{Ty-Pi} row additionally requires the
binder-edge check \(D(\sigma,\rho)\le q\).

The complete term-former list is:
\[
\begin{array}{c|c}
\text{term} & \text{principal effect}\\
\hline
x & \unit_x\\
c & \unit_c\\
\lambda^q x.t & \theta\setminus x\quad(q\le\theta(x))\\
f\,u & \phi\oplus(q\otimes\eta)\\
(t,u) & \theta\oplus\eta\\
\fst(p),\ \snd(p) & \pi\\
\mathsf{stamp}(t,e,o) & \theta\oplus\eta\oplus\omega\\
\mathsf{gate}_g(\vec d;s) &
 \unit_g\oplus\delta\oplus(1\otimes\eta)\\
\mathsf{node}_g(\vec d;t,e) &
 \unit_g\oplus\delta\oplus(1\otimes(\theta\oplus\eta))
\end{array}
\]
with exactly the typing premises in \textsc{Var}, \textsc{Const},
\textsc{Lam}, \textsc{App}, \textsc{Pair}, \textsc{Fst}, \textsc{Snd},
\textsc{Stamp}, \textsc{Gate}, and \textsc{Node}.  \textsc{Conv} is the only
rule that changes a displayed type.  \textsc{Node} is available only in
internal reducts.

\section{Reduction contexts and equality}\label{app:reduction}

Let \(C[-]\) range over one-hole term contexts:
\begin{align*}
C ::={}&[-]
 \mid \lambda^q(x:_{\sigma}A).C
 \mid C\,t\mid t\,C
 \mid(C,t)\mid(t,C)
 \mid\fst(C)\mid\snd(C)\\
 &\mid\mathsf{stamp}(C,e,o)
 \mid\mathsf{stamp}(t,C,o)
 \mid\mathsf{stamp}(t,e,C)\\
 &\mid\mathsf{gate}_g(\vec d_{<i},C,\vec d_{>i};s)
 \mid\mathsf{gate}_g(\vec d;C)\\
 &\mid\mathsf{node}_g(\vec d_{<i},C,\vec d_{>i};t,e)
 \mid\mathsf{node}_g(\vec d;C,e)
 \mid\mathsf{node}_g(\vec d;t,C).
\end{align*}
If \(t\longrightarrow t'\), then
\(C[t]\longrightarrow C[t']\) whenever both sides are well typed.

Term indices also reduce inside types.  The type contexts containing a term
hole are generated by
\begin{align*}
D ::={}&a(\vec t_{<i},C,\vec t_{>i})
 \mid\Origin_A(C,t)\mid\Origin_A(e,C)\\
 &\mid\Origin_D(e,t)
 \mid\Stamped_\sigma(D)
 \mid\AFDepFun^q(x:_{\sigma}D).B
 \mid\AFDepFun^q(x:_{\sigma}A).D\\
 &\mid\AFDepPair(x:_{\rho}D).B
 \mid\AFDepPair(x:_{\rho}A).D.
\end{align*}
The notation \(\Origin_D\) means that the payload type contains the hole.
Reduction is closed under \(D[-]\), and type contexts are also admitted in
lambda annotations.  This is every recursive position in the grammars of
\cref{sec:core}; hence the compatible closure is fully determined.

Typed definitional equality is generated by
\[
\frac{}{t\jdeq t:A}
\quad
\frac{t\jdeq u:A}{u\jdeq t:A}
\quad
\frac{t\jdeq u:A\quad u\jdeq v:A}{t\jdeq v:A}
\quad
\frac{t\longrightarrow u}{t\jdeq u:A},
\]
together with congruence for each type and term context above and conversion
between definitionally equal types.  Equality has no additional extensional,
script, quotient, proof-irrelevance, or event-erasing axiom.

\section{Reference decision procedure and implementation boundary}
\label{app:checker}

At the mathematical level, \cref{thm:decidability} induces the following
reference decision procedure for a finite proposed signature and a proposed
typing judgment:
\begin{enumerate}[label=(\arabic*),leftmargin=2.5em]
\item check the ordered signature using \textsc{Sig-Family},
\textsc{Sig-Primitive}, \textsc{Sig-Certificate},
\textsc{Sig-Attest}, and \textsc{Sig-Gate};
\item elaborate binders with explicit roles and grades;
\item synthesize types for neutral forms and check lambdas, pairs, stamps, and
other introduction forms against the supplied expected type;
\item compute principal finite-map effects bottom-up, verify
\(D(\sigma,\rho)\le q\le\theta(x)\) for an explicit lambda grade, and choose
\(q=\theta(x)\) when the grade is omitted;
\item validate every intermediate effect against \(D\);
\item normalize only when \textsc{Conv} requires comparison; and
\item for modules, compute the extension class and qualified event namespace
directly from the admitted declarations.
\end{enumerate}
Every step is deterministic after the expected judgment and ordinary
elaboration choices are fixed.  This is a mathematical algorithm
specification, not an executable or machine-verified checker.

\begin{proposition}[Declarative adequacy of the decision procedure]
The reference procedure terminates and
accepts exactly the finite signatures, contexts, types, terms, and module
classes derivable by the declarative rules.
\end{proposition}

\begin{proof}
Soundness is induction over the procedure, with one case for each
syntax-directed rule.  Completeness is induction over a declarative derivation;
principal-effect uniqueness fixes every effect annotation, while normalization
and confluence make conversion checking canonical.  Finiteness of the input and
strong normalization give termination.
\end{proof}

A concrete bidirectional judgment grammar, serialized certificate format, and
replay implementation remain open artifacts.  Such a certificate need only
record the checked signature, principal effects, module class, event owners,
and the normal forms used for conversion, but no implemented replay theorem is
claimed in this manuscript.

\clearpage
\begin{thebibliography}{99}
\small
\setlength{\itemsep}{2pt}

\bibitem{Book}
S. Krivov,
\emph{Artform Calculus},
textbook, edition of 11 October 2026.

\bibitem{Repo}
S. Krivov,
\emph{Artform Calculus: proofs in Lean},
repository, 2026,
\url{https://github.com/skrivov/artform-calculus-textbook}.

\bibitem{Lean4}
L. de Moura and S. Ullrich,
The Lean~4 theorem prover and programming language,
in \emph{Automated Deduction -- CADE 28}, LNCS 12699, Springer, 2021,
625--635.

\bibitem{Mathlib}
The mathlib Community,
The Lean mathematical library,
in \emph{Proceedings of the 9th ACM SIGPLAN International Conference on
Certified Programs and Proofs}, 2020, 367--381.

\bibitem{PatternGeometry}
S. Krivov,
\emph{Pattern Geometry in the Artform Calculus: Recognition, Generation,
Measure, Scale, Dynamics, and Evaluation of Patterns in a Typed Calculus of
Artforms},
working manuscript, 2026.

\bibitem{MartinLof}
P. Martin-L\"of,
\emph{Intuitionistic Type Theory},
Bibliopolis, 1984.

\bibitem{DybjerCwF}
P. Dybjer,
Internal type theory,
in \emph{Types for Proofs and Programs}, LNCS 1158, Springer, 1996,
120--134.

\bibitem{MTT}
D. Gratzer, G.~A. Kavvos, A. Nuyts, and L. Birkedal,
Multimodal dependent type theory,
\emph{Logical Methods in Computer Science} \textbf{17} (2021), no.~3,
article 11.

\bibitem{MTTNorm}
D. Gratzer,
Normalization for multimodal type theory,
\emph{Logical Methods in Computer Science} \textbf{22} (2026), no.~1,
article 27.

\bibitem{GradedDTT}
A. Abel, N.~A. Danielsson, and O. Eriksson,
A graded modal dependent type theory with a universe and erasure, formalized,
\emph{Proceedings of the ACM on Programming Languages} \textbf{7} (2023),
ICFP, article 220, 920--954.

\bibitem{KavvosIF}
G.~A. Kavvos,
Modalities, cohesion, and information flow,
\emph{Proceedings of the ACM on Programming Languages} \textbf{3} (2019),
POPL, article 20.

\bibitem{CoreDependency}
M. Abadi, A. Banerjee, N. Heintze, and J.~G. Riecke,
A core calculus of dependency,
in \emph{Proceedings of the 26th ACM SIGPLAN--SIGACT Symposium on Principles
of Programming Languages}, 1999, 147--160,
\url{https://doi.org/10.1145/292540.292555}.

\bibitem{Declassification}
A. Sabelfeld and D. Sands,
Dimensions and principles of declassification,
in \emph{18th IEEE Computer Security Foundations Workshop}, 2005, 255--269,
\url{https://doi.org/10.1109/CSFW.2005.15}.

\bibitem{TypeTheoryErasure}
C. Theocharis and E. Brady,
Type theory with erasure,
in \emph{11th International Conference on Formal Structures for Computation
and Deduction (FSCD 2026)}, LIPIcs 378, article 31, 1--21, 2026,
\url{https://doi.org/10.4230/LIPIcs.FSCD.2026.31}.

\bibitem{AbstractionFunctions}
H. Grodin, R. Li, and R. Harper,
Abstraction functions as types: modular verification of cost and behavior in
dependent type theory,
\emph{Proceedings of the ACM on Programming Languages} \textbf{10} (2026),
POPL, article 31, 895--922.

\bibitem{CoreProvenance}
U.~A. Acar, A. Ahmed, J. Cheney, and R. Perera,
A core calculus for provenance,
\emph{Journal of Computer Security} \textbf{21} (2013), no.~6, 919--969.

\bibitem{ModularTT}
F. Blanqui, G. Dowek, E. Grienenberger, G. Hondet, and F. Thir\'e,
A modular construction of type theories,
\emph{Logical Methods in Computer Science} \textbf{19} (2023), no.~1,
article 12.

\end{thebibliography}

\end{document}
